% Options for packages loaded elsewhere
\PassOptionsToPackage{unicode}{hyperref}
\PassOptionsToPackage{hyphens}{url}
%
\documentclass[
  12pt,
]{article}
\usepackage{lmodern}
\usepackage{amssymb,amsmath}
\usepackage{ifxetex,ifluatex}
\ifnum 0\ifxetex 1\fi\ifluatex 1\fi=0 % if pdftex
  \usepackage[T1]{fontenc}
  \usepackage[utf8]{inputenc}
  \usepackage{textcomp} % provide euro and other symbols
\else % if luatex or xetex
  \usepackage{unicode-math}
  \defaultfontfeatures{Scale=MatchLowercase}
  \defaultfontfeatures[\rmfamily]{Ligatures=TeX,Scale=1}
\fi
% Use upquote if available, for straight quotes in verbatim environments
\IfFileExists{upquote.sty}{\usepackage{upquote}}{}
\IfFileExists{microtype.sty}{% use microtype if available
  \usepackage[]{microtype}
  \UseMicrotypeSet[protrusion]{basicmath} % disable protrusion for tt fonts
}{}
\makeatletter
\@ifundefined{KOMAClassName}{% if non-KOMA class
  \IfFileExists{parskip.sty}{%
    \usepackage{parskip}
  }{% else
    \setlength{\parindent}{0pt}
    \setlength{\parskip}{6pt plus 2pt minus 1pt}}
}{% if KOMA class
  \KOMAoptions{parskip=half}}
\makeatother
\usepackage{xcolor}
\IfFileExists{xurl.sty}{\usepackage{xurl}}{} % add URL line breaks if available
\IfFileExists{bookmark.sty}{\usepackage{bookmark}}{\usepackage{hyperref}}
\hypersetup{
  pdftitle={Computing and Fitting Monotone Splines},
  pdfauthor={Jan de Leeuw},
  hidelinks,
  pdfcreator={LaTeX via pandoc}}
\urlstyle{same} % disable monospaced font for URLs
\usepackage[margin=1in]{geometry}
\usepackage{color}
\usepackage{fancyvrb}
\newcommand{\VerbBar}{|}
\newcommand{\VERB}{\Verb[commandchars=\\\{\}]}
\DefineVerbatimEnvironment{Highlighting}{Verbatim}{commandchars=\\\{\}}
% Add ',fontsize=\small' for more characters per line
\usepackage{framed}
\definecolor{shadecolor}{RGB}{248,248,248}
\newenvironment{Shaded}{\begin{snugshade}}{\end{snugshade}}
\newcommand{\AlertTok}[1]{\textcolor[rgb]{0.94,0.16,0.16}{#1}}
\newcommand{\AnnotationTok}[1]{\textcolor[rgb]{0.56,0.35,0.01}{\textbf{\textit{#1}}}}
\newcommand{\AttributeTok}[1]{\textcolor[rgb]{0.77,0.63,0.00}{#1}}
\newcommand{\BaseNTok}[1]{\textcolor[rgb]{0.00,0.00,0.81}{#1}}
\newcommand{\BuiltInTok}[1]{#1}
\newcommand{\CharTok}[1]{\textcolor[rgb]{0.31,0.60,0.02}{#1}}
\newcommand{\CommentTok}[1]{\textcolor[rgb]{0.56,0.35,0.01}{\textit{#1}}}
\newcommand{\CommentVarTok}[1]{\textcolor[rgb]{0.56,0.35,0.01}{\textbf{\textit{#1}}}}
\newcommand{\ConstantTok}[1]{\textcolor[rgb]{0.00,0.00,0.00}{#1}}
\newcommand{\ControlFlowTok}[1]{\textcolor[rgb]{0.13,0.29,0.53}{\textbf{#1}}}
\newcommand{\DataTypeTok}[1]{\textcolor[rgb]{0.13,0.29,0.53}{#1}}
\newcommand{\DecValTok}[1]{\textcolor[rgb]{0.00,0.00,0.81}{#1}}
\newcommand{\DocumentationTok}[1]{\textcolor[rgb]{0.56,0.35,0.01}{\textbf{\textit{#1}}}}
\newcommand{\ErrorTok}[1]{\textcolor[rgb]{0.64,0.00,0.00}{\textbf{#1}}}
\newcommand{\ExtensionTok}[1]{#1}
\newcommand{\FloatTok}[1]{\textcolor[rgb]{0.00,0.00,0.81}{#1}}
\newcommand{\FunctionTok}[1]{\textcolor[rgb]{0.00,0.00,0.00}{#1}}
\newcommand{\ImportTok}[1]{#1}
\newcommand{\InformationTok}[1]{\textcolor[rgb]{0.56,0.35,0.01}{\textbf{\textit{#1}}}}
\newcommand{\KeywordTok}[1]{\textcolor[rgb]{0.13,0.29,0.53}{\textbf{#1}}}
\newcommand{\NormalTok}[1]{#1}
\newcommand{\OperatorTok}[1]{\textcolor[rgb]{0.81,0.36,0.00}{\textbf{#1}}}
\newcommand{\OtherTok}[1]{\textcolor[rgb]{0.56,0.35,0.01}{#1}}
\newcommand{\PreprocessorTok}[1]{\textcolor[rgb]{0.56,0.35,0.01}{\textit{#1}}}
\newcommand{\RegionMarkerTok}[1]{#1}
\newcommand{\SpecialCharTok}[1]{\textcolor[rgb]{0.00,0.00,0.00}{#1}}
\newcommand{\SpecialStringTok}[1]{\textcolor[rgb]{0.31,0.60,0.02}{#1}}
\newcommand{\StringTok}[1]{\textcolor[rgb]{0.31,0.60,0.02}{#1}}
\newcommand{\VariableTok}[1]{\textcolor[rgb]{0.00,0.00,0.00}{#1}}
\newcommand{\VerbatimStringTok}[1]{\textcolor[rgb]{0.31,0.60,0.02}{#1}}
\newcommand{\WarningTok}[1]{\textcolor[rgb]{0.56,0.35,0.01}{\textbf{\textit{#1}}}}
\usepackage{graphicx}
\makeatletter
\def\maxwidth{\ifdim\Gin@nat@width>\linewidth\linewidth\else\Gin@nat@width\fi}
\def\maxheight{\ifdim\Gin@nat@height>\textheight\textheight\else\Gin@nat@height\fi}
\makeatother
% Scale images if necessary, so that they will not overflow the page
% margins by default, and it is still possible to overwrite the defaults
% using explicit options in \includegraphics[width, height, ...]{}
\setkeys{Gin}{width=\maxwidth,height=\maxheight,keepaspectratio}
% Set default figure placement to htbp
\makeatletter
\def\fps@figure{htbp}
\makeatother
\setlength{\emergencystretch}{3em} % prevent overfull lines
\providecommand{\tightlist}{%
  \setlength{\itemsep}{0pt}\setlength{\parskip}{0pt}}
\setcounter{secnumdepth}{5}
\newlength{\cslhangindent}
\setlength{\cslhangindent}{1.5em}
\newenvironment{cslreferences}%
  {\setlength{\parindent}{0pt}%
  \everypar{\setlength{\hangindent}{\cslhangindent}}\ignorespaces}%
  {\par}

\title{Computing and Fitting Monotone Splines}
\author{Jan de Leeuw}
\date{First created on April 17, 2017. Last update on April 13, 2020}

\begin{document}
\maketitle
\begin{abstract}
A brief introduction to spline functions and \(B\)-splines, and
specifically to monotone spline functions -- with code in R and C, and
with some applications.
\end{abstract}

{
\setcounter{tocdepth}{3}
\tableofcontents
}
Note: This is a working paper which will be expanded/updated frequently.
All suggestions for improvement are welcome. The directory
\href{http://deleeuwpdx.net/pubfolders/splines}{deleeuwpdx.net/pubfolders/splines}
has a pdf version, the complete Rmd file with all code chunks, the bib
file, and the R and C source code.

\hypertarget{introduction}{%
\section{Introduction}\label{introduction}}

To define \emph{spline functions} we first define a finite sequence of
\emph{knots} \(T=\{t_j\}\), with \(t_1\leq\cdots\leq t_p,\) and an
\emph{order} \(m\). In addition each knot \(t_j\) has a
\emph{multiplicity} \(m_j\), the number of knots equal to \(t_j\). We
suppose throughout that \(m_j\leq m\) for all \(j\).

A function \(f\) is a \emph{spline function of order} \(m\) for a knot
sequence \(\{t_j\}\) if

\begin{enumerate}
\def\labelenumi{\arabic{enumi}.}
\tightlist
\item
  \(f\) is a polynomial \(\pi_j\) of degree at most \(m-1\) on each
  half-open interval \(I_j=[t_j,t_{j+1})\) for \(j=1,\cdots,p\),
\item
  the polynomial pieces are joined in such a way that
  \(\mathcal{D}^{(s)}_-f(t_j)=\mathcal{D}^{(s)}_+f(t_j)\) for
  \(s=0,1,\cdots,m-m_j-1\) and \(j=1,2,\cdots,p\).
\end{enumerate}

Here we use \(\mathcal{D}^{(s)}_-\) and \(\mathcal{D}^{(s)}_+\) for the
left and right \(s^{th}\)-derivative operator. If \(m_j=m\) for some
\(j\), then requirement 2 is empty, if \(m_j=m-1\) then requirement 2
means \(\pi_j(t_j)=\pi_{j+1}(t_j)\), i.e.~we require continuity of \(f\)
at \(t_j\). If \(1\leq m_j<m-1\) then \(f\) must be \(m-m_j-1\) times
differentiable, and thus continuously differentiable, at \(t_j\).

In the case of simple knots (with multiplicity one) a spline function of
order one is a \emph{step function} which steps from one level to the
next at each knot. A spline of order two is piecewise linear, with the
pieces joined at the knots so that the spline function is continuous.
Order three means a piecewise quadratic function which is continuously
differentiable at the knots. And so on.

\hypertarget{basic-splines}{%
\section{Basic Splines}\label{basic-splines}}

Alternatively, a spline function of order \(m\) can be defined as a
linear combination of \emph{B-splines} (or \emph{basic splines}) of
order \(m\) on the same knot sequence. A \(B\)-spline of order \(m\) is
a spline function consisting of at most \(m\) non-zero polynomial
pieces. A \(B\)-spline \(\mathcal{B}_{j,m}\) is determined by the
\(m+1\) knots \(t_j\leq\cdots\leq t_{j+m}\), is zero outside the
interval \([t_j,t_{j+m})\), and positive in the interior of that
interval. Thus if \(t_j=t_{j+m}\) then \(\mathcal{B}_{j,m}\) is
identically zero.

For an arbitrary finite knot sequence \(t_1,\cdots,t_p\), there are
\(p-m\) \(B\)-splines to of order \(m\) to be considered, although some
may be identically zero. Each of the splines covers at most \(m\)
consecutive intervals, and at most \(m-1\) different \(B\)-splines are
non-zero at each point.

\hypertarget{boundaries}{%
\subsection{Boundaries}\label{boundaries}}

\(B\)-splines are most naturally and simply defined for doubly infinite
sequences of knots, that go to \(\pm\infty\) in both directions. In that
case we do not have to worry about boundary effects, and each
subsequence of \(m+1\) knots defines a \(B\)-spline. For splines on
finite sequences of \(p\) knots we have to decide what happens at the
boundary points.

There are \(B\)-splines for \(t_j,\cdots,t_{j+m}\) for all
\(j=1,\cdots,p-m\). This means that the first \(m-1\) and the last
\(m-1\) intervals have fewer than \(m\) splines defined on them. They
are not part of what De Boor (2001), page 94, calls the \emph{basic
interval}. For doubly infinite sequences of knots there is not need to
consider such a basic interval.

If we had \(m\) additional knots on both sides of our knot sequence we
would also have \(m\) additional \(B\)-splines for \(j=1-m,\cdots,0\)
and \(m\) additional \(B\)-splines for \(j=p-m+1,\cdots,p\). By adding
these additional knots we make sure each interval \([t_j,t_{j+1})\) for
\(j=1,\cdots,p-1\) has \(m\) \(B\)-splines associated with it. There is
stil some ambiguity on what to do at \(t_p\), but we can decide to set
the value of the spline there equal to the limit from the left, thus
making the \(B\)-spline left-continuous there.

In our software we will use the convention to define our splines on a
closed interval \([a,b]\) with \(r\) \emph{interior knots}
\(a<t_1<\cdots<t_r<b\), where interior knot \(t_j\) has multiplicity
\(m_j\). We extend this to a series of \(p=M+2m\) knots, with
\(M=\sum_{j=1}^r m_j\), by starting with \(m\) copies of \(a\),
appending \(m_j\) copies of \(t_j\) for each \(j=1,\cdots,r\), and
finishing with \(m\) copies of \(b\). Thus \(a\) and \(b\) are both
knots with multiplicity \(m\). This defines the \emph{extended
partition} (Schumaker (2007), p 116), which is just handled as any knot
sequence would normally be.

\hypertarget{normalization}{%
\subsection{Normalization}\label{normalization}}

\(B\)-splines can be defined in various ways, using piecewise
polynomials, divided differences, or recursion. The recursive
definition, first used as a defining condition by De Boor and Höllig
(1985), is the most convenient one for computational purposes, and that
is the one we use.

But first, the conditions we have mentioned only determine the
\(B\)-spline up to a normalization. There are two popular ways of
normalizing \(B\)-splines. The \(N\)-splines \(N_{j,m}\), a.k.a. the
\emph{normalized \(B\)-splines} \(j\) or order \(m\), satisfies
\begin{equation}\label{E:nsum}
\sum_{j}N_{j,m}(t)=1.
\end{equation} Note that in general this is not true for all \(t\), but
only for all \(t\) in the \emph{basic interval}.

Alternatively we can normalize to \(M\)-splines, for which
\begin{equation}\label{E:mint}
\int_{-\infty}^{+\infty}M_{j,m}(t)dt=\int_{t_j}^{t_{j+k}}M_{j,m}(t)dt=1.
\end{equation} There is the simple relationship
\begin{equation}\label{E:NM}
N_{j,m}(t)=\frac{t_{j+m}-t_j}{m}\ M_{j,m}(t).
\end{equation}

\hypertarget{recursion}{%
\subsection{Recursion}\label{recursion}}

The recursive definition, due independently to Cox (1972) for simple
knots and to De Boor (1972) in the general case, is
\begin{equation}\label{E:Mspline}
M_{j,m}(t)=\frac{t-t_j}{t_{j+m}-t_j}M_{j,m-1}(t)+\frac{t_{j+m}-t}{t_{j+m}-t_j}M_{j+1,m-1}(t),
\end{equation} or \begin{equation}\label{E:Nspline}
N_{j,m}(t)=\frac{t-t_j}{t_{m+j-1}-t_j}N_{j,m-1}(t)+\frac{t_{j+m}-t}{t_{j+m}-t_{j+1}}N_{j+1,m-1}(t).
\end{equation}

A basic result in the theory of \(B\)-splines is that the different
\(B\)-splines are linearly independent and form a basis for the linear
space of spline functions (of a given order and knot sequence).

\hypertarget{computations}{%
\section{Computations}\label{computations}}

Before introducing our new C code we review some the approaches we have
used in the past. This will also give us the opportunity to make some
comparisons.

\hypertarget{low-order-splines}{%
\subsection{Low Order Splines}\label{low-order-splines}}

The R code in \texttt{lowSpline.R} has three functions to compute
splines of order one, two, and three. It does not acknowledge any
boundary values, so only looks at \(B\)-splines on an interval spanned
by interior knots. The formulas we use are

\[
N_{j,1}(x)=\begin{cases}1&\text{ if }t_j\leq x<t_{j+1},\\0&\text{ otherwise}\end{cases}.
\]

\[
N_{j,2}(x)=\begin{cases}\frac{x-t_j}{t_{j+1}-t_j}&\text{ if }t_j\leq x<t_{j+1},\\\frac{t_{j+2}-x}{t_{j+2}-t_{j+1}}&\text{ if }t_{j+1}\leq x<t_{j+2},\\0&\text{ otherwise}\end{cases}.
\]

\[
N_{j,3}(x)=\begin{cases}\frac{(x-t_j)^2}{(t_{j+1}-t_j)(t_{j+2}-t_j)}&\text{ if }t_j\leq x<t_{j+1},\\
\frac{(x-t_j)(t_{j+2}-x)}{(t_{j+2}-t_j)(t_{j+2}-t_{j+1})}+\frac{(x-t_{j+1})(t_{j+3}-x)}{(t_{j+3}-t_{j+1})(t_{j+2}-t_{j+1})}&\text{ if }t_{j+1}\leq x<t_{j+2},\\
\frac{(t_{j+3}-x)^2}{(t_{j+3}-t_{j+1})(t_{j+3}-t_{j+2})}&\text{ if }t_{j+2}\leq x<t_{j+3},\\
0&\text{ otherwise}\end{cases}.
\]

In the example of Ramsay (1988) the knots are 0.0, 0.3, 0.5, 0.6, and
1.0 (in the example 0 and 1 are endpoints of the interval, but we'll
just treat them as interior knots in a longer sequence). Also note that
Ramsay computes \(M\)-splines, while we compute \(N\)-splines.

We start with the \(p-m=5-2=3\) \(B\)-splines of order 2. The basic
interval is \([0.3,0.6]\).

\begin{center}\includegraphics{splines_files/figure-latex/lowsecond-1} \end{center}

Figure 1: Piecewise Linear Splines with Simple Knots

There are only two \(B\)-splines of order 3 on these knots, and the
basic interval is empty.

\begin{center}\includegraphics{splines_files/figure-latex/lowthird-1} \end{center}

Figure 2: Piecewise Quadratic Splines with Simple Knots

The programs also work for multiple knots. Consider the example from De
Boor (2001), page 92. The knots are 0, 1, 1, 3, 4, 6, 6, 6, and the
order is three. The basic interval is \([1,6]\).

\begin{center}\includegraphics{splines_files/figure-latex/lowthirdmult-1} \end{center}

Figure 3: Piecewise Quadratic Splines with Multiple Knots

In De Boor (2001), p 89, we find another example in which there are only
two distinct knots, an infinite sequence of zeroes, followed by an
infinite sequence of ones. In this case there are only \(m\) different
\(B\)-splines, restriction to \([0,1]\) of the familiar polynomials \[
B_{j,m-1}(x)=\binom{m-1}{j}x^{m-1-j}(1-x)^{j}
\] for \(j=0,\cdots,m-1\).

\begin{center}\includegraphics{splines_files/figure-latex/lowthirdmultb-1} \end{center}

Figure 4: Quadratic Bernstein Basis

\hypertarget{using-gsl}{%
\subsection{Using GSL}\label{using-gsl}}

The GNU Scientific Library (Galassi et al. (2016)) has \(B\)-spline
code. The function \texttt{gslSpline()} in R calls the compiled
\texttt{gslSpline.c}, which is linked with the relevant code from
\texttt{libgsl.dylib}. We use the Ramsay example again. The GSL
implementation automatically adds the extra boundary knots for the
extended partition, which makes the basic interval \([0,1]\).

\begin{Shaded}
\begin{Highlighting}[]
\NormalTok{knots \textless{}{-}}\StringTok{ }\KeywordTok{c}\NormalTok{(}\DecValTok{0}\NormalTok{,.}\DecValTok{3}\NormalTok{,.}\DecValTok{5}\NormalTok{,.}\DecValTok{6}\NormalTok{,}\DecValTok{1}\NormalTok{)}
\NormalTok{order \textless{}{-}}\StringTok{ }\DecValTok{3}
\NormalTok{x\textless{}{-}}\KeywordTok{seq}\NormalTok{(}\DecValTok{0}\NormalTok{,}\DecValTok{1}\NormalTok{,}\DataTypeTok{length =} \DecValTok{1001}\NormalTok{)}
\NormalTok{h \textless{}{-}}\StringTok{ }\KeywordTok{matrix}\NormalTok{ (}\DecValTok{0}\NormalTok{, }\DecValTok{1001}\NormalTok{, }\DecValTok{6}\NormalTok{)}
\ControlFlowTok{for}\NormalTok{ (i }\ControlFlowTok{in} \DecValTok{1}\OperatorTok{:}\DecValTok{1001}\NormalTok{)}
\NormalTok{  h[i,] \textless{}{-}}\StringTok{ }\KeywordTok{gslSpline}\NormalTok{ (x[i], order, knots)}
\end{Highlighting}
\end{Shaded}

\begin{center}\includegraphics{splines_files/figure-latex/with_gsl_plot-1} \end{center}

Figure 5: Piecewise Quadratic Splines using GSL

\hypertarget{using-recursion}{%
\subsection{Using Recursion}\label{using-recursion}}

We have previously published spline function code, using an R interface
to C code, in J De Leeuw (2015). That code translated the Fortran code
in an unpublished note by Sinha (n.d.) to C. There are some limitations
associated with this implementation. First, it is limited to extended
partitions with simple inner knots. Second, the function to compute
\(B\)-spline values recursively calls itself, using the basic relation
\(\eqref{E:Nspline}\), which is computationally not necessarily the best
way to go.

\begin{Shaded}
\begin{Highlighting}[]
\NormalTok{innerknots \textless{}{-}}\StringTok{ }\KeywordTok{c}\NormalTok{( }\FloatTok{.3}\NormalTok{, }\FloatTok{.5}\NormalTok{, }\FloatTok{.6}\NormalTok{)}
\NormalTok{degree \textless{}{-}}\StringTok{ }\DecValTok{2}
\NormalTok{lowknot \textless{}{-}}\StringTok{ }\DecValTok{0}
\NormalTok{highknot \textless{}{-}}\StringTok{ }\DecValTok{1}
\NormalTok{x\textless{}{-}}\KeywordTok{seq}\NormalTok{(}\DecValTok{0}\NormalTok{,}\DecValTok{1}\NormalTok{,}\DataTypeTok{length =} \DecValTok{1001}\NormalTok{)}
\NormalTok{h \textless{}{-}}\StringTok{ }\KeywordTok{sinhaBasis}\NormalTok{ (x, degree, innerknots, lowknot, highknot)}
\end{Highlighting}
\end{Shaded}

\begin{center}\includegraphics{splines_files/figure-latex/with_sinha_plot-1} \end{center}

Figure 6: Piecewise Quadratic Splinesusing Recursion

\hypertarget{de-boor}{%
\subsection{De Boor}\label{de-boor}}

In this paper we implement the basic BSPLVB algorithm from De Boor
(2001), page 111, for normalized \(B\)-splines. There are two auxilary
routines, one to create the extended partition, and one that uses
bisection to locate the knot interval in which a particular value is
located (Schumaker (2007), p 191). The function \texttt{bsplineBasis()}
takes an arbitrary knot sequence. It can be combined with
\texttt{extendPartition()}, which uses inner knots and boundary points
to create the extended partion.

\hypertarget{illustrations}{%
\subsection{Illustrations}\label{illustrations}}

For our example, which is the same as the one from figure 1 in Ramsay
(1988), we choose \(a=0\), \(b=1\), with simple interior knots 0.3, 0.5,
0.6. First the step functions, which have order 1.

\begin{Shaded}
\begin{Highlighting}[]
\NormalTok{innerknots \textless{}{-}}\StringTok{ }\KeywordTok{c}\NormalTok{(.}\DecValTok{3}\NormalTok{,.}\DecValTok{5}\NormalTok{,.}\DecValTok{6}\NormalTok{)}
\NormalTok{multiplicities \textless{}{-}}\StringTok{ }\KeywordTok{c}\NormalTok{(}\DecValTok{1}\NormalTok{,}\DecValTok{1}\NormalTok{,}\DecValTok{1}\NormalTok{)}
\NormalTok{order \textless{}{-}}\StringTok{ }\DecValTok{1}
\NormalTok{lowend \textless{}{-}}\StringTok{ }\DecValTok{0}
\NormalTok{highend \textless{}{-}}\StringTok{ }\DecValTok{1}
\NormalTok{x \textless{}{-}}\StringTok{ }\KeywordTok{seq}\NormalTok{ (}\FloatTok{1e{-}6}\NormalTok{, }\DecValTok{1}\FloatTok{{-}1e{-}6}\NormalTok{, }\DataTypeTok{length =} \DecValTok{1000}\NormalTok{)}
\NormalTok{knots \textless{}{-}}\StringTok{ }\KeywordTok{extendPartition}\NormalTok{ (innerknots, multiplicities, order, lowend, highend)}\OperatorTok{$}\NormalTok{knots}
\NormalTok{h \textless{}{-}}\StringTok{ }\KeywordTok{bsplineBasis}\NormalTok{ (x, knots, order)}
\NormalTok{k \textless{}{-}}\StringTok{ }\KeywordTok{ncol}\NormalTok{ (h)}
\KeywordTok{par}\NormalTok{ (}\DataTypeTok{mfrow=}\KeywordTok{c}\NormalTok{(}\DecValTok{2}\NormalTok{,}\DecValTok{2}\NormalTok{))}
\ControlFlowTok{for}\NormalTok{ (j }\ControlFlowTok{in} \DecValTok{1}\OperatorTok{:}\NormalTok{k) \{ }
\NormalTok{  ylab \textless{}{-}}\StringTok{ }\KeywordTok{paste}\NormalTok{(}\StringTok{"B{-}spline"}\NormalTok{, }\KeywordTok{formatC}\NormalTok{(j, }\DataTypeTok{digits =} \DecValTok{1}\NormalTok{, }\DataTypeTok{width =} \DecValTok{2}\NormalTok{, }\DataTypeTok{format =} \StringTok{"d"}\NormalTok{))}
  \KeywordTok{plot}\NormalTok{ (x, h[, j], }\DataTypeTok{type=}\StringTok{"l"}\NormalTok{, }\DataTypeTok{col =} \StringTok{"RED"}\NormalTok{, }\DataTypeTok{lwd =} \DecValTok{3}\NormalTok{, }\DataTypeTok{ylab =}\NormalTok{ ylab, }\DataTypeTok{ylim =} \KeywordTok{c}\NormalTok{(}\DecValTok{0}\NormalTok{,}\DecValTok{1}\NormalTok{))}
\NormalTok{\}}
\end{Highlighting}
\end{Shaded}

\begin{center}\includegraphics{splines_files/figure-latex/order1mult1-1} \end{center}

Figure 7: Zero Degree Splines with Simple Knots

Now the hat functions, which have order 2, again with simple knots.

\begin{Shaded}
\begin{Highlighting}[]
\NormalTok{multiplicities \textless{}{-}}\StringTok{ }\KeywordTok{c}\NormalTok{(}\DecValTok{1}\NormalTok{,}\DecValTok{1}\NormalTok{,}\DecValTok{1}\NormalTok{)}
\NormalTok{order \textless{}{-}}\StringTok{ }\DecValTok{2}
\NormalTok{knots \textless{}{-}}\StringTok{ }\KeywordTok{extendPartition}\NormalTok{ (innerknots, multiplicities, order, lowend, highend)}\OperatorTok{$}\NormalTok{knots}
\NormalTok{h \textless{}{-}}\StringTok{ }\KeywordTok{bsplineBasis}\NormalTok{ (x, knots, order)}
\NormalTok{k \textless{}{-}}\StringTok{ }\KeywordTok{ncol}\NormalTok{ (h)}
\KeywordTok{par}\NormalTok{ (}\DataTypeTok{mfrow=}\KeywordTok{c}\NormalTok{(}\DecValTok{2}\NormalTok{,}\DecValTok{3}\NormalTok{))}
\ControlFlowTok{for}\NormalTok{ (j }\ControlFlowTok{in} \DecValTok{1}\OperatorTok{:}\NormalTok{k) \{ }
\NormalTok{  ylab \textless{}{-}}\StringTok{ }\KeywordTok{paste}\NormalTok{(}\StringTok{"B{-}spline"}\NormalTok{, }\KeywordTok{formatC}\NormalTok{(j, }\DataTypeTok{digits =} \DecValTok{1}\NormalTok{, }\DataTypeTok{width =} \DecValTok{2}\NormalTok{, }\DataTypeTok{format =} \StringTok{"d"}\NormalTok{))}
  \KeywordTok{plot}\NormalTok{ (x, h[, j], }\DataTypeTok{type=}\StringTok{"l"}\NormalTok{, }\DataTypeTok{col =} \StringTok{"RED"}\NormalTok{, }\DataTypeTok{lwd =} \DecValTok{3}\NormalTok{, }\DataTypeTok{ylab =}\NormalTok{ ylab, }\DataTypeTok{ylim =} \KeywordTok{c}\NormalTok{(}\DecValTok{0}\NormalTok{,}\DecValTok{1}\NormalTok{))}
\NormalTok{\}}
\end{Highlighting}
\end{Shaded}

\begin{center}\includegraphics{splines_files/figure-latex/order2mult1-1} \end{center}

Figure 8: Piecewise Linear Splines with Simple Knots

Next piecewise quadratics, with simple knots, which implies continuous
differentiability at the knots. This are the N-splines corresponding
with the M-splines in figure 1 of Ramsay (1988).

\begin{Shaded}
\begin{Highlighting}[]
\NormalTok{multiplicities \textless{}{-}}\StringTok{ }\KeywordTok{c}\NormalTok{(}\DecValTok{1}\NormalTok{,}\DecValTok{1}\NormalTok{,}\DecValTok{1}\NormalTok{)}
\NormalTok{order \textless{}{-}}\StringTok{ }\DecValTok{3}
\NormalTok{knots \textless{}{-}}\StringTok{ }\KeywordTok{extendPartition}\NormalTok{ (innerknots, multiplicities, order, lowend, highend)}\OperatorTok{$}\NormalTok{knots}
\NormalTok{h \textless{}{-}}\StringTok{ }\KeywordTok{bsplineBasis}\NormalTok{ (x, knots, order)}
\NormalTok{k \textless{}{-}}\StringTok{ }\KeywordTok{ncol}\NormalTok{ (h)}
\KeywordTok{par}\NormalTok{ (}\DataTypeTok{mfrow=}\KeywordTok{c}\NormalTok{(}\DecValTok{2}\NormalTok{,}\DecValTok{3}\NormalTok{))}
\ControlFlowTok{for}\NormalTok{ (j }\ControlFlowTok{in} \DecValTok{1}\OperatorTok{:}\NormalTok{k) \{ }
\NormalTok{  ylab \textless{}{-}}\StringTok{ }\KeywordTok{paste}\NormalTok{(}\StringTok{"B{-}spline"}\NormalTok{, }\KeywordTok{formatC}\NormalTok{(j, }\DataTypeTok{digits =} \DecValTok{1}\NormalTok{, }\DataTypeTok{width =} \DecValTok{2}\NormalTok{, }\DataTypeTok{format =} \StringTok{"d"}\NormalTok{))}
  \KeywordTok{plot}\NormalTok{ (x, h[, j], }\DataTypeTok{type=}\StringTok{"l"}\NormalTok{, }\DataTypeTok{col =} \StringTok{"RED"}\NormalTok{, }\DataTypeTok{lwd =} \DecValTok{3}\NormalTok{, }\DataTypeTok{ylab =}\NormalTok{ ylab, }\DataTypeTok{ylim =} \KeywordTok{c}\NormalTok{(}\DecValTok{0}\NormalTok{,}\DecValTok{1}\NormalTok{))}
\NormalTok{\}}
\end{Highlighting}
\end{Shaded}

\begin{center}\includegraphics{splines_files/figure-latex/order3mult1-1} \end{center}

Figure 9: Piecewise Quadratic Splines with Simple Knots

If we change the multiplicities to 1, 2, 3, then we lose some of the
smoothness.

\begin{Shaded}
\begin{Highlighting}[]
\NormalTok{multiplicities \textless{}{-}}\StringTok{ }\KeywordTok{c}\NormalTok{(}\DecValTok{1}\NormalTok{,}\DecValTok{2}\NormalTok{,}\DecValTok{3}\NormalTok{)}
\NormalTok{order \textless{}{-}}\StringTok{ }\DecValTok{3}
\NormalTok{knots \textless{}{-}}\StringTok{ }\KeywordTok{extendPartition}\NormalTok{ (innerknots, multiplicities, order, lowend, highend)}\OperatorTok{$}\NormalTok{knots}
\NormalTok{h \textless{}{-}}\StringTok{ }\KeywordTok{bsplineBasis}\NormalTok{ (x, knots, order)}
\NormalTok{k \textless{}{-}}\StringTok{ }\KeywordTok{ncol}\NormalTok{ (h)}
\KeywordTok{par}\NormalTok{ (}\DataTypeTok{mfrow=}\KeywordTok{c}\NormalTok{(}\DecValTok{3}\NormalTok{,}\DecValTok{3}\NormalTok{))}
\ControlFlowTok{for}\NormalTok{ (j }\ControlFlowTok{in} \DecValTok{1}\OperatorTok{:}\NormalTok{k) \{ }
\NormalTok{  ylab \textless{}{-}}\StringTok{ }\KeywordTok{paste}\NormalTok{(}\StringTok{"B{-}spline"}\NormalTok{, }\KeywordTok{formatC}\NormalTok{(j, }\DataTypeTok{digits =} \DecValTok{1}\NormalTok{, }\DataTypeTok{width =} \DecValTok{2}\NormalTok{, }\DataTypeTok{format =} \StringTok{"d"}\NormalTok{))}
  \KeywordTok{plot}\NormalTok{ (x, h[, j], }\DataTypeTok{type=}\StringTok{"l"}\NormalTok{, }\DataTypeTok{col =} \StringTok{"RED"}\NormalTok{, }\DataTypeTok{lwd =} \DecValTok{3}\NormalTok{, }\DataTypeTok{ylab =}\NormalTok{ ylab, }\DataTypeTok{ylim =} \KeywordTok{c}\NormalTok{(}\DecValTok{0}\NormalTok{,}\DecValTok{1}\NormalTok{))}
\NormalTok{\}}
\end{Highlighting}
\end{Shaded}

\begin{center}\includegraphics{splines_files/figure-latex/order3mult2-1} \end{center}

Figure 10: Piecewise Quadratic Splines with Multiple Knots

\hypertarget{monotone-splines}{%
\section{Monotone Splines}\label{monotone-splines}}

\hypertarget{i-splines}{%
\subsection{I-splines}\label{i-splines}}

There are several ways to restrict splines to be monotone increasing.
Since \(B\)-splines are non-negative, the definite integral of a
\(B\)-spline of order \(m\) from the beginning of the interval to a
value \(x\) in the interval is an increasing spline of order \(m+1\).
Integrated \(B\)-splines are known as \emph{I-splines} (Ramsay (1988)).
Non-negative linear combinations \(I\)-splines can be used as a basis
for the convex cone of increasing splines. Note, however, that if we use
an extended partition, then all \(I\)-splines start at value zero and
end at value one, which means their convex combinations are the splines
that are also probability distributions on the interval. To get a basis
for the increasing splines we need to add the constant function to the
\(I\)-splines and allow it to enter the linear combination with either
sign.

\hypertarget{low-order-i-splines}{%
\subsubsection{Low Order I-splines}\label{low-order-i-splines}}

Straightforward integration, and using \(\eqref{E:NM}\), gives some
explicit formulas. If we integrate the step functions we get the
piecewise linear \(I\)-splines. \[
\int_{-\infty}^x M_{j,1}(t)dt=\begin{cases}
0&\text{ if }x\leq t_j,\\
\frac{x-t_j}{t_{j+1}-t_j}&\text{ if }t_j\leq x\leq t_{j+1},\\
1&\text{ if }x\geq t_{j+1}.
\end{cases}
\] And if we integrate the piecewise linear \(B\)-splines of order 1 we
get piecewise quadratic \(I\)-splines. \[
\int_{-\infty}^x M_{j,2}(t)dt=\begin{cases}
0&\text{ if }x\leq t_j,\\
\frac{(x-t_j)^2}{(t_{j+1}-t_j)(t_{j+2}- t_j)}&\text{ if }t_j\leq x\leq t_{j+1},\\
\frac{x-t_j}{t_{j+2}-t_j}+\frac{(x-t_{j+1})(t_{j+2}-x)}{(t_{j+2}-t_j)(t_{j+2}-t_{j+1})}
&\text{ if }t_{j+1}\leq x\leq t_{j+2},\\
1&\text{ if }x\geq t_{j+2}.
\end{cases}
\] Both sets of \(I\)-splines are plotted in the next two figures, using
R functions from \texttt{lowSpline.R}.

\begin{center}\includegraphics{splines_files/figure-latex/ilowfirst-1} \end{center}

Figure 11: Monotone Piecewise Linear Splines with Simple Knots

\begin{center}\includegraphics{splines_files/figure-latex/ilowsecond-1} \end{center}

Figure 12: Monotone Piecewise Quadratic Splines with Simple Knots

\hypertarget{general-case}{%
\subsubsection{General Case}\label{general-case}}

Integrals of I-splines are most economically computed by using the
formula first given by Gaffney (1976). If \(\ell\) is defined by
\(t_{j+\ell-1}\leq x<t_{j+\ell}\) then \[
\int_{x_j}^x M_{j,m}(t)dt=\frac{1}{m}\sum_{r=0}^{ \ell-1}(x-x_{j+r})M_{j+r,m-r}(x)
\] It is somewhat simpler, however, to use lemma 2.1 of De Boor, Lyche,
and Schumaker (1976). This says \[
\int_a^xM_{j,m}(t)dt=\sum_{\ell\geq j}N_{\ell,m+1}(x)-\sum_{\ell\geq j}N_{\ell,m+1}(a),
\] If we specialize this to \(I\)-splines, we find , as in De Boor
(1976), formula 4.11, \[
\int_{-\infty}^x M_{j,m}(t)dt=\sum_{\ell=j}^{j+r}N_{\ell,m+1}(x)
\] for \(x\leq t_{j+r+1}\). This shows that \(I\)-splines can be
computed by using cumulative sums of \(B\)-spline values.

Note that using the integration definition does not give a natural way
to define increasing splines of degree 1, i.e.~increasing step
functions. There is no such problem with the cumulative sum approach.

\hypertarget{increasing-coefficients}{%
\subsection{Increasing Coefficients}\label{increasing-coefficients}}

As we know, a spline is a linear combination of \(B\)-splines. The
formula for the derivative of a spline, for example in De Boor (2001), p
116, shows that a spline is increasing if the coefficients of the linear
combination of \(B\)-splines are increasing. Thus we can fit an
increasing spline by restricting the coefficients of the linear
combination to be increasing, again using the \(B\)-spline basis.

It turns out this is in fact identical to using \(I\)-splines. If the
\(B\)-spline values at \(n\) points are in an \(n\times r\) matrix
\(H\), then increasing coefficients \(\beta\) are of the form
\(\beta=S\alpha+\gamma e_r\), where \(S\) is lower-diagonal with all
elements on and below the diagonal equal to one, where \(\alpha\geq 0\),
where \(e_r\) has all elements equal to one, and where \(\gamma\) can be
of any sign. So \(H\beta=(HS)\alpha+\gamma e_n\). The matrix \(Z=HS\) is
easily and cheaply found in R by \(1-t(apply(h, 1, cumsum))\).

\hypertarget{increasing-values}{%
\subsection{Increasing Values}\label{increasing-values}}

Finally, we can simply require that the \(n\) elements of \(H\beta\) are
increasing. This is a less restrictive requirement, because it allows
for the possibility that the spline is decresing between data values. It
has the rather serious disadvantage, however, that it does its
computations in \(n\)-dimensional space, and not in \(r\)-dimensional
space, where \(r=M+m\), which is usually much smaller than \(n\).
Software for the increasing-value restrictions has been written by J. De
Leeuw (2015). In this paper, however, we prefer the \texttt{cumsum()}
approach. It is less general, but considerably more efficient.

\hypertarget{illustrations-1}{%
\subsection{Illustrations}\label{illustrations-1}}

We use the same Ramsay example as before, but now cumulatively. First we
integrate step functions with simple knots, which have order 1, using
\texttt{isplineBasis()}. The corresponding I-splines are piecewise
linear with order 2.

\begin{center}\includegraphics{splines_files/figure-latex/Iorder1mult1-1} \end{center}

Figure 23:

Now we integrate the hat functions, which have order 2, again with
simple knots, to find piecewise quadratic I-splines of order 3. These
are the functions in the example of Ramsay (1988).

\begin{center}\includegraphics{splines_files/figure-latex/Iorder2mult1-1} \end{center}

Figure 13: Monotone Piecewise Linear Splines with Simple Knots

Finally, we change the multiplicities to 1, 2, 3, and compute the
corresponding piecewise quadratic I-splines.

\begin{center}\includegraphics{splines_files/figure-latex/Iorder3mult2-1} \end{center}

Figure 15: Monotone Piecewise Quadratic Splines with Multiple Knots

\hypertarget{time-series-example}{%
\section{Time Series Example}\label{time-series-example}}

Our first example smoothes a time series by fitting a spline. We use the
number of births in New York from 1946 to 1959 (on an unknown scale),
from Rob Hyndman's time series archive at
\url{http://robjhyndman.com/tsdldata/data/nybirths.dat}.

\hypertarget{b-splines}{%
\subsection{B-splines}\label{b-splines}}

First we fit \(B\)-splines of order 3. The basis matrix uses \(x\) equal
to \(1:168\), with inner knots 12, 24, 36, 48, 60, 72, 84, 96, 108, 120,
132, 144, 156, and interval \([1,168]\).

\begin{Shaded}
\begin{Highlighting}[]
\NormalTok{innerknots \textless{}{-}}\StringTok{ }\DecValTok{12} \OperatorTok{*}\StringTok{ }\DecValTok{1}\OperatorTok{:}\DecValTok{13}
\NormalTok{multiplicities \textless{}{-}}\StringTok{ }\KeywordTok{rep}\NormalTok{(}\DecValTok{1}\NormalTok{,}\DecValTok{13}\NormalTok{)}
\NormalTok{lowend \textless{}{-}}\StringTok{ }\DecValTok{1}
\NormalTok{highend \textless{}{-}}\StringTok{ }\DecValTok{168}
\NormalTok{order \textless{}{-}}\StringTok{ }\DecValTok{3}
\NormalTok{x \textless{}{-}}\StringTok{ }\DecValTok{1}\OperatorTok{:}\DecValTok{168}
\NormalTok{knots \textless{}{-}}\StringTok{ }\KeywordTok{extendPartition}\NormalTok{ (innerknots, multiplicities, order, lowend, highend)}\OperatorTok{$}\NormalTok{knots}
\NormalTok{h \textless{}{-}}\StringTok{ }\KeywordTok{bsplineBasis}\NormalTok{ (x, knots, order)}
\NormalTok{u \textless{}{-}}\StringTok{ }\KeywordTok{lm.fit}\NormalTok{(h, births)}
\NormalTok{res \textless{}{-}}\StringTok{ }\KeywordTok{sum}\NormalTok{ ((births }\OperatorTok{{-}}\StringTok{ }\NormalTok{h}\OperatorTok{\%*\%}\NormalTok{u}\OperatorTok{$}\NormalTok{coefficients)}\OperatorTok{\^{}}\DecValTok{2}\NormalTok{)}
\end{Highlighting}
\end{Shaded}

\begin{center}\includegraphics{splines_files/figure-latex/birth-bspline-plot-1} \end{center}

Figure 16: Monotone Piecewise Quadratic Splines with Simple Knots

The residual sum of squares is 229.3835417745.

\hypertarget{i-splines-1}{%
\subsection{I-splines}\label{i-splines-1}}

We now fit the \(I\)-spline using the \(B\)-spline basis. Compute
\(Z=HS\) using \texttt{cumsum()}, and then \(\overline y\) and
\(\overline Z\) by centering (substracting the column means). The
formula is \[
\min_{\alpha\geq 0,\gamma}\mathbf{SSQ}\ (y-Z\alpha-\gamma e_n)=\min_{\alpha\geq 0}\mathbf{SSQ}\ (\overline y-\overline Z\alpha).
\] We use \texttt{pnnls()} from Wang, Lawson, and Hanson (2015).

\begin{Shaded}
\begin{Highlighting}[]
\NormalTok{knots \textless{}{-}}\StringTok{ }\KeywordTok{extendPartition}\NormalTok{ (innerknots, multiplicities, order, lowend, highend)}\OperatorTok{$}\NormalTok{knots}
\NormalTok{h \textless{}{-}}\StringTok{ }\KeywordTok{isplineBasis}\NormalTok{ (x, knots, order)}
\NormalTok{g \textless{}{-}}\StringTok{ }\KeywordTok{cbind}\NormalTok{ (}\DecValTok{1}\NormalTok{, h[,}\OperatorTok{{-}}\DecValTok{1}\NormalTok{])}
\NormalTok{u \textless{}{-}}\StringTok{ }\KeywordTok{pnnls}\NormalTok{ (g, births, }\DecValTok{1}\NormalTok{)}\OperatorTok{$}\NormalTok{x}
\NormalTok{v \textless{}{-}}\StringTok{ }\NormalTok{g}\OperatorTok{\%*\%}\NormalTok{u}
\end{Highlighting}
\end{Shaded}

\begin{center}\includegraphics{splines_files/figure-latex/births_isplines_plot-1} \end{center}

Figure 17: Monotone Piecewise Linear Splines with Simple Knots

The residual sum of squares is 288.4054982424.

\hypertarget{b-splines-with-monotone-weights}{%
\subsection{B-Splines with monotone
weights}\label{b-splines-with-monotone-weights}}

Just to make sure, we also solve the problem \[
\min_{\beta_1\leq\beta_2\leq\cdots\leq\beta_p}\mathbf{SSQ}(y-X\beta),
\] which should give the same solution, and the same loss function
value, because it is just another way to fit \(I\)-splines. We use the
\texttt{lsi()} function from Wang, Lawson, and Hanson (2015).

\begin{Shaded}
\begin{Highlighting}[]
\NormalTok{knots \textless{}{-}}\StringTok{ }\KeywordTok{extendPartition}\NormalTok{ (innerknots, multiplicities, order, lowend, highend)}\OperatorTok{$}\NormalTok{knots}
\NormalTok{h \textless{}{-}}\StringTok{ }\KeywordTok{bsplineBasis}\NormalTok{ (x, knots, order)}
\NormalTok{nb \textless{}{-}}\StringTok{ }\KeywordTok{ncol}\NormalTok{ (h)}
\NormalTok{d \textless{}{-}}\StringTok{ }\KeywordTok{matrix}\NormalTok{(}\DecValTok{0}\NormalTok{,nb}\DecValTok{{-}1}\NormalTok{,nb)}
\KeywordTok{diag}\NormalTok{(d)=}\OperatorTok{{-}}\DecValTok{1}
\NormalTok{d[}\KeywordTok{outer}\NormalTok{(}\DecValTok{1}\OperatorTok{:}\NormalTok{(nb}\DecValTok{{-}1}\NormalTok{),}\DecValTok{1}\OperatorTok{:}\NormalTok{nb,}\ControlFlowTok{function}\NormalTok{(i,j) (j }\OperatorTok{{-}}\StringTok{ }\NormalTok{i) }\OperatorTok{==}\StringTok{ }\DecValTok{1}\NormalTok{)]\textless{}{-}}\DecValTok{1}
\NormalTok{u\textless{}{-}}\KeywordTok{lsi}\NormalTok{(h,births,}\DataTypeTok{e=}\NormalTok{d,}\DataTypeTok{f=}\KeywordTok{rep}\NormalTok{(}\DecValTok{0}\NormalTok{,nb}\DecValTok{{-}1}\NormalTok{))}
\NormalTok{v \textless{}{-}}\StringTok{ }\NormalTok{h }\OperatorTok{\%*\%}\StringTok{ }\NormalTok{u}
\end{Highlighting}
\end{Shaded}

\begin{center}\includegraphics{splines_files/figure-latex/births_isplines_values_plot-1} \end{center}

Figure 18: Monotone Piecewise Quadratic Splines with Simple Knots

The residual sum of squares is 288.4054982424, indeed the same as
before.

\hypertarget{b-splines-with-monotone-values}{%
\subsection{B-Splines with monotone
values}\label{b-splines-with-monotone-values}}

Finally we solve \[
\min_{x_1'\beta\leq\cdots\leq x_n'\beta} \mathbf{SSQ}\ (y-X\beta)
\] using \texttt{mregnnM()} from J De Leeuw (2015), which solves the
dual problem using \texttt{nnls()} from Wang, Lawson, and Hanson (2015).

\begin{Shaded}
\begin{Highlighting}[]
\NormalTok{knots \textless{}{-}}\StringTok{ }\KeywordTok{extendPartition}\NormalTok{ (innerknots, multiplicities, order, lowend, highend)}\OperatorTok{$}\NormalTok{knots}
\NormalTok{h \textless{}{-}}\StringTok{ }\KeywordTok{bsplineBasis}\NormalTok{ (x, knots, order)}
\NormalTok{u \textless{}{-}}\StringTok{ }\KeywordTok{mregnnM}\NormalTok{(h, births)}
\end{Highlighting}
\end{Shaded}

\begin{center}\includegraphics{splines_files/figure-latex/births_bsplines_monotonic_plot-1} \end{center}

Figure 19: Monotone Piecewise Quadratic Splines with Multiple Knots

The residual sum of squares is 288.3210359867, which is indeed smaller
than the \(I\)-splines value, although only very slightly so.

\hypertarget{regression-example}{%
\section{Regression Example}\label{regression-example}}

We also analyze a regression example, using the Neumann data that have
been analyzed previously in Gifi (1990), pages 370-376, and in De Leeuw
and Mair (2009),pages 16-17. The predictors are temperature and
pressure, the outcome variable is density. We use a step function
monotone spline for temperature and a piecewise quadratic monotone
spline for pressure. The lest squares problem is to fit \[
\mathbf{SSQ}(y-(\gamma e_n+H_1\alpha_1+H_2\alpha_2)),
\] where \(H_1\) and \(H_2\) are \(I\)-spline bases and \(\alpha_1\) and
\(\alpha_2\) are non-negative. We do not transform the outcome variable
density, to keep things relatively simple.

\begin{Shaded}
\begin{Highlighting}[]
\KeywordTok{data}\NormalTok{(neumann, }\DataTypeTok{package =} \StringTok{"homals"}\NormalTok{)}
\NormalTok{knots1 \textless{}{-}}\StringTok{ }\KeywordTok{c}\NormalTok{(}\DecValTok{0}\NormalTok{, }\DecValTok{20}\NormalTok{, }\DecValTok{40}\NormalTok{, }\DecValTok{60}\NormalTok{, }\DecValTok{80}\NormalTok{, }\DecValTok{100}\NormalTok{, }\DecValTok{120}\NormalTok{, }\DecValTok{140}\NormalTok{)}
\NormalTok{order1 \textless{}{-}}\StringTok{ }\DecValTok{1}
\NormalTok{knots2 \textless{}{-}}\StringTok{ }\KeywordTok{c}\NormalTok{(}\DecValTok{0}\NormalTok{,}\DecValTok{0}\NormalTok{,}\DecValTok{0}\NormalTok{,}\DecValTok{100}\NormalTok{,}\DecValTok{200}\NormalTok{,}\DecValTok{300}\NormalTok{,}\DecValTok{400}\NormalTok{,}\DecValTok{500}\NormalTok{,}\DecValTok{600}\NormalTok{,}\DecValTok{600}\NormalTok{,}\DecValTok{600}\NormalTok{)}
\NormalTok{order2 \textless{}{-}}\StringTok{ }\DecValTok{3}
\NormalTok{h1 \textless{}{-}}\StringTok{ }\KeywordTok{isplineBasis}\NormalTok{(}\DecValTok{200}\OperatorTok{{-}}\NormalTok{neumann[,}\DecValTok{1}\NormalTok{], knots1, order1)}
\NormalTok{h2 \textless{}{-}}\StringTok{ }\KeywordTok{isplineBasis}\NormalTok{(neumann[,}\DecValTok{2}\NormalTok{], knots2, order2)}
\NormalTok{g \textless{}{-}}\StringTok{ }\KeywordTok{cbind}\NormalTok{ (}\DecValTok{1}\NormalTok{, h1, h2)}
\NormalTok{u \textless{}{-}}\StringTok{ }\KeywordTok{pnnls}\NormalTok{ (g, neumann[,}\DecValTok{3}\NormalTok{], }\DecValTok{1}\NormalTok{)}\OperatorTok{$}\NormalTok{x}
\NormalTok{u1 \textless{}{-}}\StringTok{ }\NormalTok{u[}\DecValTok{1}\OperatorTok{+}\DecValTok{1}\OperatorTok{:}\KeywordTok{ncol}\NormalTok{(h1)]}
\NormalTok{u2 \textless{}{-}}\StringTok{ }\NormalTok{u[}\DecValTok{1}\OperatorTok{+}\KeywordTok{ncol}\NormalTok{(h1)}\OperatorTok{+}\DecValTok{1}\OperatorTok{:}\KeywordTok{ncol}\NormalTok{(h2)]}
\end{Highlighting}
\end{Shaded}

The next two plots give the transformations of the predictors.

\begin{center}\includegraphics{splines_files/figure-latex/neumann_plot1-1} \end{center}

Figure 20: Regression Example: Piecewise Constant Temperature

\begin{center}\includegraphics{splines_files/figure-latex/neumann_plot2-1} \end{center}

Figure 21: Regression Example: Piecewise Quadratic Pressure

And finally, we give the residual plot, observed density versus
predicted density.

\begin{center}\includegraphics{splines_files/figure-latex/neumann_plot3-1} \end{center}

Figure 22: Regression Example: Fit

\hypertarget{appendix-code}{%
\section{Appendix: Code}\label{appendix-code}}

\hypertarget{r-code}{%
\subsection{R code}\label{r-code}}

\hypertarget{lowspline.r}{%
\subsubsection{lowSpline.R}\label{lowspline.r}}

\begin{Shaded}
\begin{Highlighting}[]
\NormalTok{ZbSpline \textless{}{-}}\StringTok{ }\ControlFlowTok{function}\NormalTok{(x, knots, }\DataTypeTok{k =} \DecValTok{1}\NormalTok{) \{}
\NormalTok{  ZbSplineSingle \textless{}{-}}\StringTok{ }\ControlFlowTok{function}\NormalTok{(x, knots, }\DataTypeTok{k =} \DecValTok{1}\NormalTok{) \{}
\NormalTok{    k0 \textless{}{-}}\StringTok{ }\NormalTok{knots[k]}
\NormalTok{    k1 \textless{}{-}}\StringTok{ }\NormalTok{knots[k }\OperatorTok{+}\StringTok{ }\DecValTok{1}\NormalTok{]}
    \ControlFlowTok{if}\NormalTok{ ((x }\OperatorTok{\textgreater{}}\StringTok{ }\NormalTok{k0) }\OperatorTok{\&\&}\StringTok{ }\NormalTok{(x }\OperatorTok{\textless{}=}\StringTok{ }\NormalTok{k1)) \{}
      \KeywordTok{return}\NormalTok{(}\DecValTok{1}\NormalTok{)}
\NormalTok{    \}}
    \KeywordTok{return}\NormalTok{(}\DecValTok{0}\NormalTok{)}
\NormalTok{  \}}
  \KeywordTok{return}\NormalTok{(}\KeywordTok{sapply}\NormalTok{(x, }\ControlFlowTok{function}\NormalTok{(z)}
    \KeywordTok{ZbSplineSingle}\NormalTok{(z, knots, k)))}
\NormalTok{\}}

\NormalTok{LbSpline \textless{}{-}}\StringTok{ }\ControlFlowTok{function}\NormalTok{(x, knots, }\DataTypeTok{k =} \DecValTok{1}\NormalTok{) \{}
\NormalTok{  LbSplineSingle \textless{}{-}}\StringTok{ }\ControlFlowTok{function}\NormalTok{(x, knots, }\DataTypeTok{k =} \DecValTok{1}\NormalTok{) \{}
\NormalTok{    k0 \textless{}{-}}\StringTok{ }\NormalTok{knots[k]}
\NormalTok{    k1 \textless{}{-}}\StringTok{ }\NormalTok{knots[k }\OperatorTok{+}\StringTok{ }\DecValTok{1}\NormalTok{]}
\NormalTok{    k2 \textless{}{-}}\StringTok{ }\NormalTok{knots[k }\OperatorTok{+}\StringTok{ }\DecValTok{2}\NormalTok{]}
\NormalTok{    f1 \textless{}{-}}\StringTok{ }\ControlFlowTok{function}\NormalTok{(x)}
\NormalTok{      (x }\OperatorTok{{-}}\StringTok{ }\NormalTok{k0) }\OperatorTok{/}\StringTok{ }\NormalTok{(k1 }\OperatorTok{{-}}\StringTok{ }\NormalTok{k0)}
\NormalTok{    f2 \textless{}{-}}\StringTok{ }\ControlFlowTok{function}\NormalTok{(x)}
\NormalTok{      (k2 }\OperatorTok{{-}}\StringTok{ }\NormalTok{x) }\OperatorTok{/}\StringTok{ }\NormalTok{(k2 }\OperatorTok{{-}}\StringTok{ }\NormalTok{k1)}
    \ControlFlowTok{if}\NormalTok{ ((x }\OperatorTok{\textgreater{}}\StringTok{ }\NormalTok{k0) }\OperatorTok{\&\&}\StringTok{ }\NormalTok{(x }\OperatorTok{\textless{}=}\StringTok{ }\NormalTok{k1)) \{}
      \KeywordTok{return}\NormalTok{(}\KeywordTok{f1}\NormalTok{(x))}
\NormalTok{    \}}
    \ControlFlowTok{if}\NormalTok{ ((x }\OperatorTok{\textgreater{}}\StringTok{ }\NormalTok{k1) }\OperatorTok{\&\&}\StringTok{ }\NormalTok{(x }\OperatorTok{\textless{}=}\StringTok{ }\NormalTok{k2)) \{}
      \KeywordTok{return}\NormalTok{(}\KeywordTok{f2}\NormalTok{(x))}
\NormalTok{    \}}
    \KeywordTok{return}\NormalTok{(}\DecValTok{0}\NormalTok{)}
\NormalTok{  \}}
  \KeywordTok{return}\NormalTok{(}\KeywordTok{sapply}\NormalTok{(x, }\ControlFlowTok{function}\NormalTok{(z)}
    \KeywordTok{LbSplineSingle}\NormalTok{(z, knots, k)))}
\NormalTok{\}}

\NormalTok{QbSpline \textless{}{-}}\StringTok{ }\ControlFlowTok{function}\NormalTok{(x, knots, }\DataTypeTok{k =} \DecValTok{1}\NormalTok{) \{}
\NormalTok{  QbSplineSingle \textless{}{-}}\StringTok{ }\ControlFlowTok{function}\NormalTok{(x, knots, }\DataTypeTok{k =} \DecValTok{1}\NormalTok{) \{}
\NormalTok{    k0 \textless{}{-}}\StringTok{ }\NormalTok{knots[k]}
\NormalTok{    k1 \textless{}{-}}\StringTok{ }\NormalTok{knots[k }\OperatorTok{+}\StringTok{ }\DecValTok{1}\NormalTok{]}
\NormalTok{    k2 \textless{}{-}}\StringTok{ }\NormalTok{knots[k }\OperatorTok{+}\StringTok{ }\DecValTok{2}\NormalTok{]}
\NormalTok{    k3 \textless{}{-}}\StringTok{ }\NormalTok{knots[k }\OperatorTok{+}\StringTok{ }\DecValTok{3}\NormalTok{]}
\NormalTok{    f1 \textless{}{-}}\StringTok{ }\ControlFlowTok{function}\NormalTok{(x)}
\NormalTok{      ((x }\OperatorTok{{-}}\StringTok{ }\NormalTok{k0) }\OperatorTok{\^{}}\StringTok{ }\DecValTok{2}\NormalTok{) }\OperatorTok{/}\StringTok{ }\NormalTok{((k2 }\OperatorTok{{-}}\StringTok{ }\NormalTok{k0) }\OperatorTok{*}\StringTok{ }\NormalTok{(k1 }\OperatorTok{{-}}\StringTok{ }\NormalTok{k0))}
\NormalTok{    f2 \textless{}{-}}\StringTok{ }\ControlFlowTok{function}\NormalTok{(x) \{}
\NormalTok{      term1 \textless{}{-}}\StringTok{ }\NormalTok{((x }\OperatorTok{{-}}\StringTok{ }\NormalTok{k0) }\OperatorTok{*}\StringTok{ }\NormalTok{(k2 }\OperatorTok{{-}}\StringTok{ }\NormalTok{x)) }\OperatorTok{/}\StringTok{ }\NormalTok{((k2 }\OperatorTok{{-}}\StringTok{ }\NormalTok{k0) }\OperatorTok{*}\StringTok{ }\NormalTok{(k2 }\OperatorTok{{-}}\StringTok{ }\NormalTok{k1))}
\NormalTok{      term2 \textless{}{-}}\StringTok{ }\NormalTok{((x }\OperatorTok{{-}}\StringTok{ }\NormalTok{k1) }\OperatorTok{*}\StringTok{ }\NormalTok{(k3 }\OperatorTok{{-}}\StringTok{ }\NormalTok{x)) }\OperatorTok{/}\StringTok{ }\NormalTok{((k3 }\OperatorTok{{-}}\StringTok{ }\NormalTok{k1) }\OperatorTok{*}\StringTok{ }\NormalTok{(k2 }\OperatorTok{{-}}\StringTok{ }\NormalTok{k1))}
      \KeywordTok{return}\NormalTok{(term1 }\OperatorTok{+}\StringTok{ }\NormalTok{term2)}
\NormalTok{    \}}
\NormalTok{    f3 \textless{}{-}}\StringTok{ }\ControlFlowTok{function}\NormalTok{(x)}
\NormalTok{      ((k3 }\OperatorTok{{-}}\StringTok{ }\NormalTok{x) }\OperatorTok{\^{}}\StringTok{ }\DecValTok{2}\NormalTok{) }\OperatorTok{/}\StringTok{ }\NormalTok{((k3 }\OperatorTok{{-}}\StringTok{ }\NormalTok{k1) }\OperatorTok{*}\StringTok{ }\NormalTok{(k3 }\OperatorTok{{-}}\StringTok{ }\NormalTok{k2))}
    \ControlFlowTok{if}\NormalTok{ ((x }\OperatorTok{\textgreater{}}\StringTok{ }\NormalTok{k0) }\OperatorTok{\&\&}\StringTok{ }\NormalTok{(x }\OperatorTok{\textless{}=}\StringTok{ }\NormalTok{k1)) \{}
      \KeywordTok{return}\NormalTok{(}\KeywordTok{f1}\NormalTok{(x))}
\NormalTok{    \}}
    \ControlFlowTok{if}\NormalTok{ ((x }\OperatorTok{\textgreater{}}\StringTok{ }\NormalTok{k1) }\OperatorTok{\&\&}\StringTok{ }\NormalTok{(x }\OperatorTok{\textless{}=}\StringTok{ }\NormalTok{k2)) \{}
      \KeywordTok{return}\NormalTok{(}\KeywordTok{f2}\NormalTok{(x))}
\NormalTok{    \}}
    \ControlFlowTok{if}\NormalTok{ ((x }\OperatorTok{\textgreater{}}\StringTok{ }\NormalTok{k2) }\OperatorTok{\&\&}\StringTok{ }\NormalTok{(x }\OperatorTok{\textless{}=}\StringTok{ }\NormalTok{k3)) \{}
      \KeywordTok{return}\NormalTok{(}\KeywordTok{f3}\NormalTok{(x))}
\NormalTok{    \}}
    \KeywordTok{return}\NormalTok{(}\DecValTok{0}\NormalTok{)}
\NormalTok{  \}}
  \KeywordTok{return}\NormalTok{(}\KeywordTok{sapply}\NormalTok{(x, }\ControlFlowTok{function}\NormalTok{(z)}
    \KeywordTok{QbSplineSingle}\NormalTok{(z, knots, k)))}
\NormalTok{\}}

\NormalTok{IZbSpline \textless{}{-}}\StringTok{ }\ControlFlowTok{function}\NormalTok{(x, knots, }\DataTypeTok{k =} \DecValTok{1}\NormalTok{) \{}
\NormalTok{  IZbSplineSingle \textless{}{-}}\StringTok{ }\ControlFlowTok{function}\NormalTok{(x, knots, }\DataTypeTok{k =} \DecValTok{1}\NormalTok{) \{}
\NormalTok{    k0 \textless{}{-}}\StringTok{ }\NormalTok{knots[k]}
\NormalTok{    k1 \textless{}{-}}\StringTok{ }\NormalTok{knots[k }\OperatorTok{+}\StringTok{ }\DecValTok{1}\NormalTok{]}
    \ControlFlowTok{if}\NormalTok{ (x }\OperatorTok{\textless{}=}\StringTok{ }\NormalTok{k0) \{}
      \KeywordTok{return}\NormalTok{ (}\DecValTok{0}\NormalTok{)}
\NormalTok{    \}}
    \ControlFlowTok{if}\NormalTok{ ((x }\OperatorTok{\textgreater{}}\StringTok{ }\NormalTok{k0) }\OperatorTok{\&\&}\StringTok{ }\NormalTok{(x }\OperatorTok{\textless{}=}\StringTok{ }\NormalTok{k1)) \{}
      \KeywordTok{return}\NormalTok{((x }\OperatorTok{{-}}\StringTok{ }\NormalTok{k0) }\OperatorTok{/}\StringTok{ }\NormalTok{(k1 }\OperatorTok{{-}}\StringTok{ }\NormalTok{k0))}
\NormalTok{    \}}
    \ControlFlowTok{if}\NormalTok{ (x }\OperatorTok{\textgreater{}}\StringTok{ }\NormalTok{k1) \{}
      \KeywordTok{return}\NormalTok{ (}\DecValTok{1}\NormalTok{)}
\NormalTok{    \}}
\NormalTok{  \}}
  \KeywordTok{return}\NormalTok{(}\KeywordTok{sapply}\NormalTok{(x, }\ControlFlowTok{function}\NormalTok{(z)}
    \KeywordTok{IZbSplineSingle}\NormalTok{(z, knots, k)))}
\NormalTok{\}}

\NormalTok{ILbSpline \textless{}{-}}\StringTok{ }\ControlFlowTok{function}\NormalTok{(x, knots, }\DataTypeTok{k =} \DecValTok{1}\NormalTok{) \{}
\NormalTok{  ILbSplineSingle \textless{}{-}}\StringTok{ }\ControlFlowTok{function}\NormalTok{(x, knots, }\DataTypeTok{k =} \DecValTok{1}\NormalTok{) \{}
\NormalTok{    k0 \textless{}{-}}\StringTok{ }\NormalTok{knots[k]}
\NormalTok{    k1 \textless{}{-}}\StringTok{ }\NormalTok{knots[k }\OperatorTok{+}\StringTok{ }\DecValTok{1}\NormalTok{]}
\NormalTok{    k2 \textless{}{-}}\StringTok{ }\NormalTok{knots[k }\OperatorTok{+}\StringTok{ }\DecValTok{2}\NormalTok{]}
\NormalTok{    f1 \textless{}{-}}\StringTok{ }\ControlFlowTok{function}\NormalTok{(x)}
\NormalTok{      ((x }\OperatorTok{{-}}\StringTok{ }\NormalTok{k0) }\OperatorTok{\^{}}\StringTok{ }\DecValTok{2}\NormalTok{) }\OperatorTok{/}\StringTok{ }\NormalTok{((k1 }\OperatorTok{{-}}\StringTok{ }\NormalTok{k0) }\OperatorTok{*}\StringTok{ }\NormalTok{(k2 }\OperatorTok{{-}}\StringTok{ }\NormalTok{k0))}
\NormalTok{    f2 \textless{}{-}}
\StringTok{      }\ControlFlowTok{function}\NormalTok{(x)}
\NormalTok{        ((x }\OperatorTok{{-}}\StringTok{ }\NormalTok{k0) }\OperatorTok{/}\StringTok{ }\NormalTok{(k2 }\OperatorTok{{-}}\StringTok{ }\NormalTok{k0)) }\OperatorTok{+}\StringTok{ }\NormalTok{((x }\OperatorTok{{-}}\StringTok{ }\NormalTok{k1) }\OperatorTok{*}\StringTok{ }\NormalTok{(k2 }\OperatorTok{{-}}\StringTok{ }\NormalTok{x)) }\OperatorTok{/}\StringTok{ }\NormalTok{((k2 }\OperatorTok{{-}}\StringTok{ }\NormalTok{k0) }\OperatorTok{*}\StringTok{ }\NormalTok{(k2 }\OperatorTok{{-}}\StringTok{ }\NormalTok{k1))}
    \ControlFlowTok{if}\NormalTok{ (x }\OperatorTok{\textless{}=}\StringTok{ }\NormalTok{k0) \{}
      \KeywordTok{return}\NormalTok{ (}\DecValTok{0}\NormalTok{)}
\NormalTok{    \}}
    \ControlFlowTok{if}\NormalTok{ ((x }\OperatorTok{\textgreater{}}\StringTok{ }\NormalTok{k0) }\OperatorTok{\&\&}\StringTok{ }\NormalTok{(x }\OperatorTok{\textless{}=}\StringTok{ }\NormalTok{k1)) \{}
      \KeywordTok{return}\NormalTok{(}\KeywordTok{f1}\NormalTok{(x))}
\NormalTok{    \}}
    \ControlFlowTok{if}\NormalTok{ ((x }\OperatorTok{\textgreater{}}\StringTok{ }\NormalTok{k1) }\OperatorTok{\&\&}\StringTok{ }\NormalTok{(x }\OperatorTok{\textless{}=}\StringTok{ }\NormalTok{k2)) \{}
      \KeywordTok{return}\NormalTok{(}\KeywordTok{f2}\NormalTok{(x))}
\NormalTok{    \}}
    \ControlFlowTok{if}\NormalTok{ (x }\OperatorTok{\textgreater{}}\StringTok{ }\NormalTok{k2) \{}
      \KeywordTok{return}\NormalTok{ (}\DecValTok{1}\NormalTok{)}
\NormalTok{    \}}
    \KeywordTok{return}\NormalTok{(}\DecValTok{0}\NormalTok{)}
\NormalTok{  \}}
  \KeywordTok{return}\NormalTok{(}\KeywordTok{sapply}\NormalTok{(x, }\ControlFlowTok{function}\NormalTok{(z)}
    \KeywordTok{ILbSplineSingle}\NormalTok{(z, knots, k)))}
\NormalTok{\}}
\end{Highlighting}
\end{Shaded}

\hypertarget{gslspline.r}{%
\subsubsection{gslSpline.R}\label{gslspline.r}}

\begin{Shaded}
\begin{Highlighting}[]
\KeywordTok{dyn.load}\NormalTok{(}\StringTok{"gslSpline.so"}\NormalTok{)}

\NormalTok{gslSpline \textless{}{-}}\StringTok{ }\ControlFlowTok{function}\NormalTok{ (x, k, br) \{}
\NormalTok{  nbr \textless{}{-}}\StringTok{ }\KeywordTok{length}\NormalTok{ (br)}
\NormalTok{  nrs \textless{}{-}}\StringTok{ }\NormalTok{k }\OperatorTok{+}\StringTok{ }\NormalTok{nbr }\OperatorTok{{-}}\StringTok{ }\DecValTok{2}
\NormalTok{  res \textless{}{-}}
\StringTok{    }\KeywordTok{.C}\NormalTok{(}
      \StringTok{"BSPLINE"}\NormalTok{,}
      \KeywordTok{as.double}\NormalTok{ (x),}
      \KeywordTok{as.integer}\NormalTok{ (k),}
      \KeywordTok{as.integer}\NormalTok{ (nbr),}
      \KeywordTok{as.double}\NormalTok{ (br),}
      \DataTypeTok{results =} \KeywordTok{as.double}\NormalTok{ (}\KeywordTok{rep}\NormalTok{(}\FloatTok{0.0}\NormalTok{, nrs))}
\NormalTok{    )}
  \KeywordTok{return}\NormalTok{ (res}\OperatorTok{$}\NormalTok{results)}
\NormalTok{\}}
\end{Highlighting}
\end{Shaded}

\hypertarget{sinhaspline.r}{%
\subsubsection{sinhaSpline.R}\label{sinhaspline.r}}

\begin{Shaded}
\begin{Highlighting}[]
\KeywordTok{dyn.load}\NormalTok{(}\StringTok{"sinhaSpline.so"}\NormalTok{)}

\NormalTok{sinhaBasis \textless{}{-}}
\StringTok{  }\ControlFlowTok{function}\NormalTok{ (x, degree, innerknots, }\DataTypeTok{lowknot =} \KeywordTok{min}\NormalTok{(x,innerknots), }\DataTypeTok{highknot =} \KeywordTok{max}\NormalTok{(x,innerknots)) \{}
\NormalTok{    innerknots \textless{}{-}}\StringTok{ }\KeywordTok{unique}\NormalTok{ (}\KeywordTok{sort}\NormalTok{ (innerknots))}
\NormalTok{    knots \textless{}{-}}
\StringTok{      }\KeywordTok{c}\NormalTok{(}\KeywordTok{rep}\NormalTok{(lowknot, degree }\OperatorTok{+}\StringTok{ }\DecValTok{1}\NormalTok{), innerknots, }\KeywordTok{rep}\NormalTok{(highknot, degree }\OperatorTok{+}\StringTok{ }\DecValTok{1}\NormalTok{))}
\NormalTok{    n \textless{}{-}}\StringTok{ }\KeywordTok{length}\NormalTok{ (x)}
\NormalTok{    m \textless{}{-}}\StringTok{ }\KeywordTok{length}\NormalTok{ (innerknots) }\OperatorTok{+}\StringTok{ }\DecValTok{2} \OperatorTok{*}\StringTok{ }\NormalTok{(degree }\OperatorTok{+}\StringTok{ }\DecValTok{1}\NormalTok{)}
\NormalTok{    nf \textless{}{-}}\StringTok{ }\KeywordTok{length}\NormalTok{ (innerknots) }\OperatorTok{+}\StringTok{ }\NormalTok{degree }\OperatorTok{+}\StringTok{ }\DecValTok{1}
\NormalTok{    basis \textless{}{-}}\StringTok{ }\KeywordTok{rep}\NormalTok{ (}\DecValTok{0}\NormalTok{,  n }\OperatorTok{*}\StringTok{ }\NormalTok{nf)}
\NormalTok{    res \textless{}{-}}\StringTok{ }\KeywordTok{.C}\NormalTok{(}
      \StringTok{"sinhaBasis"}\NormalTok{, }\DataTypeTok{d =} \KeywordTok{as.integer}\NormalTok{(degree),}
      \DataTypeTok{n =} \KeywordTok{as.integer}\NormalTok{(n), }\DataTypeTok{m =} \KeywordTok{as.integer}\NormalTok{ (m), }\DataTypeTok{x =} \KeywordTok{as.double}\NormalTok{ (x), }\DataTypeTok{knots =} \KeywordTok{as.double}\NormalTok{ (knots), }\DataTypeTok{basis =} \KeywordTok{as.double}\NormalTok{(basis)}
\NormalTok{    )}
\NormalTok{    basis \textless{}{-}}\StringTok{ }\KeywordTok{matrix}\NormalTok{ (res}\OperatorTok{$}\NormalTok{basis, n, nf)}
\NormalTok{    basis \textless{}{-}}\StringTok{ }\NormalTok{basis[,}\KeywordTok{which}\NormalTok{(}\KeywordTok{colSums}\NormalTok{(basis) }\OperatorTok{\textgreater{}}\StringTok{ }\DecValTok{0}\NormalTok{)]}
    \KeywordTok{return}\NormalTok{ (basis)}
\NormalTok{  \}}
\end{Highlighting}
\end{Shaded}

\hypertarget{deboor.r}{%
\subsubsection{deboor.R}\label{deboor.r}}

\begin{Shaded}
\begin{Highlighting}[]
\KeywordTok{dyn.load}\NormalTok{(}\StringTok{"deboor.so"}\NormalTok{)}

\NormalTok{checkIncreasing \textless{}{-}}\StringTok{ }\ControlFlowTok{function}\NormalTok{ (innerknots, lowend, highend) \{}
\NormalTok{  h \textless{}{-}}\StringTok{ }\KeywordTok{.C}\NormalTok{(}
    \StringTok{"checkIncreasing"}\NormalTok{,}
    \KeywordTok{as.double}\NormalTok{ (innerknots),}
    \KeywordTok{as.double}\NormalTok{ (lowend),}
    \KeywordTok{as.double}\NormalTok{ (highend),}
    \KeywordTok{as.integer}\NormalTok{ (}\KeywordTok{length}\NormalTok{ (innerknots)),}
    \DataTypeTok{fail =} \KeywordTok{as.integer}\NormalTok{ (}\DecValTok{0}\NormalTok{)}
\NormalTok{  )}
  \KeywordTok{return}\NormalTok{ (h}\OperatorTok{$}\NormalTok{fail)}
\NormalTok{\}}

\NormalTok{extendPartition \textless{}{-}}
\StringTok{  }\ControlFlowTok{function}\NormalTok{ (innerknots,}
\NormalTok{            multiplicities,}
\NormalTok{            order,}
\NormalTok{            lowend,}
\NormalTok{            highend) \{}
\NormalTok{    ninner \textless{}{-}}\StringTok{ }\KeywordTok{length}\NormalTok{ (innerknots)}
\NormalTok{    kk \textless{}{-}}\StringTok{ }\KeywordTok{sum}\NormalTok{(multiplicities)}
\NormalTok{    nextended \textless{}{-}}\StringTok{ }\NormalTok{kk }\OperatorTok{+}\StringTok{ }\DecValTok{2} \OperatorTok{*}\StringTok{ }\NormalTok{order}
    \ControlFlowTok{if}\NormalTok{ (}\KeywordTok{max}\NormalTok{ (multiplicities) }\OperatorTok{\textgreater{}}\StringTok{ }\NormalTok{order)}
      \KeywordTok{stop}\NormalTok{ (}\StringTok{"multiplicities too high"}\NormalTok{)}
    \ControlFlowTok{if}\NormalTok{ (}\KeywordTok{min}\NormalTok{ (multiplicities) }\OperatorTok{\textless{}}\StringTok{ }\DecValTok{1}\NormalTok{)}
      \KeywordTok{stop}\NormalTok{ (}\StringTok{"multiplicities too low"}\NormalTok{)}
    \ControlFlowTok{if}\NormalTok{ (}\KeywordTok{checkIncreasing}\NormalTok{ (innerknots, lowend, highend))}
      \KeywordTok{stop}\NormalTok{ (}\StringTok{"knot sequence not increasing"}\NormalTok{)}
\NormalTok{    h \textless{}{-}}
\StringTok{      }\KeywordTok{.C}\NormalTok{(}
        \StringTok{"extendPartition"}\NormalTok{,}
        \KeywordTok{as.double}\NormalTok{ (innerknots),}
        \KeywordTok{as.integer}\NormalTok{ (multiplicities),}
        \KeywordTok{as.integer}\NormalTok{ (order),}
        \KeywordTok{as.integer}\NormalTok{ (ninner),}
        \KeywordTok{as.double}\NormalTok{ (lowend),}
        \KeywordTok{as.double}\NormalTok{ (highend),}
        \DataTypeTok{knots =} \KeywordTok{as.double}\NormalTok{ (}\KeywordTok{rep}\NormalTok{ (}\DecValTok{0}\NormalTok{, nextended))}
\NormalTok{      )}
    \KeywordTok{return}\NormalTok{ (h)}
\NormalTok{  \}}

\NormalTok{bisect \textless{}{-}}
\StringTok{  }\ControlFlowTok{function}\NormalTok{ (x,}
\NormalTok{            knots,}
            \DataTypeTok{lowindex =} \DecValTok{1}\NormalTok{,}
            \DataTypeTok{highindex =} \KeywordTok{length}\NormalTok{ (knots)) \{}
\NormalTok{    h \textless{}{-}}\StringTok{ }\KeywordTok{.C}\NormalTok{(}
      \StringTok{"bisect"}\NormalTok{,}
      \KeywordTok{as.double}\NormalTok{ (x),}
      \KeywordTok{as.double}\NormalTok{ (knots),}
      \KeywordTok{as.integer}\NormalTok{ (lowindex),}
      \KeywordTok{as.integer}\NormalTok{ (highindex),}
      \DataTypeTok{index =} \KeywordTok{as.integer}\NormalTok{ (}\DecValTok{0}\NormalTok{)}
\NormalTok{    )}
    \KeywordTok{return}\NormalTok{ (h}\OperatorTok{$}\NormalTok{index)}
\NormalTok{  \}}

\NormalTok{bsplines \textless{}{-}}\StringTok{ }\ControlFlowTok{function}\NormalTok{ (x, knots, order) \{}
  \ControlFlowTok{if}\NormalTok{ ((x }\OperatorTok{\textgreater{}}\StringTok{ }\NormalTok{knots[}\KeywordTok{length}\NormalTok{(knots)]) }\OperatorTok{||}\StringTok{ }\NormalTok{(x }\OperatorTok{\textless{}}\StringTok{ }\NormalTok{knots[}\DecValTok{1}\NormalTok{]))}
    \KeywordTok{stop}\NormalTok{ (}\StringTok{"argument out of range"}\NormalTok{)}
\NormalTok{  h \textless{}{-}}\StringTok{ }\KeywordTok{.C}\NormalTok{(}
    \StringTok{"bsplines"}\NormalTok{,}
    \KeywordTok{as.double}\NormalTok{ (x),}
    \KeywordTok{as.double}\NormalTok{ (knots),}
    \KeywordTok{as.integer}\NormalTok{ (order),}
    \KeywordTok{as.integer}\NormalTok{ (}\KeywordTok{length}\NormalTok{ (knots)),}
    \DataTypeTok{index =} \KeywordTok{as.integer}\NormalTok{ (}\DecValTok{0}\NormalTok{),}
    \DataTypeTok{q =} \KeywordTok{as.double}\NormalTok{ (}\KeywordTok{rep}\NormalTok{(}\DecValTok{0}\NormalTok{, order))}
\NormalTok{  )}
  \KeywordTok{return}\NormalTok{ (}\KeywordTok{list}\NormalTok{ (}\DataTypeTok{q =}\NormalTok{ h}\OperatorTok{$}\NormalTok{q, }\DataTypeTok{index =}\NormalTok{ h}\OperatorTok{$}\NormalTok{ind))}
\NormalTok{\}}


\NormalTok{bsplineBasis \textless{}{-}}\StringTok{ }\ControlFlowTok{function}\NormalTok{ (x, knots, order) \{}
\NormalTok{  n \textless{}{-}}\StringTok{ }\KeywordTok{length}\NormalTok{ (x)}
\NormalTok{  k \textless{}{-}}\StringTok{ }\KeywordTok{length}\NormalTok{ (knots)}
\NormalTok{  m \textless{}{-}}\StringTok{ }\NormalTok{k }\OperatorTok{{-}}\StringTok{ }\NormalTok{order}
\NormalTok{  result \textless{}{-}}\StringTok{ }\KeywordTok{rep}\NormalTok{ (}\DecValTok{0}\NormalTok{, n }\OperatorTok{*}\StringTok{ }\NormalTok{m)}
\NormalTok{  h \textless{}{-}}\StringTok{ }\KeywordTok{.C}\NormalTok{(}
    \StringTok{"bsplineBasis"}\NormalTok{,}
    \KeywordTok{as.double}\NormalTok{ (x),}
    \KeywordTok{as.double}\NormalTok{ (knots),}
    \KeywordTok{as.integer}\NormalTok{ (order),}
    \KeywordTok{as.integer}\NormalTok{ (k),}
    \KeywordTok{as.integer}\NormalTok{ (n),}
    \DataTypeTok{result =} \KeywordTok{as.double}\NormalTok{ (result)}
\NormalTok{  )}
  \KeywordTok{return}\NormalTok{ (}\KeywordTok{matrix}\NormalTok{ (h}\OperatorTok{$}\NormalTok{result, n, m))}
\NormalTok{\}}

\NormalTok{isplineBasis \textless{}{-}}\StringTok{  }\ControlFlowTok{function}\NormalTok{ (x, knots, order) \{}
\NormalTok{  n \textless{}{-}}\StringTok{ }\KeywordTok{length}\NormalTok{ (x)}
\NormalTok{  k \textless{}{-}}\StringTok{ }\KeywordTok{length}\NormalTok{ (knots)}
\NormalTok{  m \textless{}{-}}\StringTok{ }\NormalTok{k }\OperatorTok{{-}}\StringTok{ }\NormalTok{order}
\NormalTok{  result \textless{}{-}}\StringTok{ }\KeywordTok{rep}\NormalTok{ (}\DecValTok{0}\NormalTok{, n }\OperatorTok{*}\StringTok{ }\NormalTok{m)}
\NormalTok{  h \textless{}{-}}\StringTok{ }\KeywordTok{.C}\NormalTok{(}
    \StringTok{"isplineBasis"}\NormalTok{,}
    \KeywordTok{as.double}\NormalTok{ (x),}
    \KeywordTok{as.double}\NormalTok{ (knots),}
    \KeywordTok{as.integer}\NormalTok{ (order),}
    \KeywordTok{as.integer}\NormalTok{ (k),}
    \KeywordTok{as.integer}\NormalTok{ (n),}
    \DataTypeTok{result =} \KeywordTok{as.double}\NormalTok{ (result)}
\NormalTok{  )}
  \KeywordTok{return}\NormalTok{ (}\KeywordTok{matrix}\NormalTok{ (h}\OperatorTok{$}\NormalTok{result, n, m))}
\NormalTok{\}}
\end{Highlighting}
\end{Shaded}

\hypertarget{c-code}{%
\subsection{C code}\label{c-code}}

\hypertarget{gslspline.c}{%
\subsubsection{gslSpline.c}\label{gslspline.c}}

\begin{Shaded}
\begin{Highlighting}[]
\PreprocessorTok{\#include }\ImportTok{\textless{}gsl/gsl\_bspline.h\textgreater{}}

\DataTypeTok{void}\NormalTok{ BSPLINE(}\DataTypeTok{double}\NormalTok{ *x, }\DataTypeTok{int}\NormalTok{ *order, }\DataTypeTok{int}\NormalTok{ *nbreak, }\DataTypeTok{double}\NormalTok{ *brpnts,}
             \DataTypeTok{double}\NormalTok{ *results) \{}
\NormalTok{    gsl\_bspline\_workspace *my\_workspace =}
\NormalTok{        gsl\_bspline\_alloc((}\DataTypeTok{size\_t}\NormalTok{)*order, (}\DataTypeTok{size\_t}\NormalTok{)*nbreak);}
    \DataTypeTok{size\_t}\NormalTok{ ncoefs = gsl\_bspline\_ncoeffs(my\_workspace);}
\NormalTok{    gsl\_vector *values = gsl\_vector\_calloc(ncoefs);}
\NormalTok{    gsl\_vector *breaks = gsl\_vector\_calloc((}\DataTypeTok{size\_t}\NormalTok{)*nbreak);}
    \ControlFlowTok{for}\NormalTok{ (}\DataTypeTok{int}\NormalTok{ i = }\DecValTok{0}\NormalTok{; i \textless{} *nbreak; i++)}
\NormalTok{        gsl\_vector\_set(breaks, (}\DataTypeTok{size\_t}\NormalTok{)i, brpnts[i]);}
\NormalTok{    (}\DataTypeTok{void}\NormalTok{)gsl\_bspline\_knots(breaks, my\_workspace);}
\NormalTok{    (}\DataTypeTok{void}\NormalTok{)gsl\_bspline\_eval(*x, values, my\_workspace);}
    \ControlFlowTok{for}\NormalTok{ (}\DataTypeTok{int}\NormalTok{ i = }\DecValTok{0}\NormalTok{; i \textless{} ncoefs; i++) results[i] = (values{-}\textgreater{}data)[i];}
\NormalTok{    gsl\_bspline\_free(my\_workspace);}
\NormalTok{    gsl\_vector\_free(values);}
\NormalTok{    gsl\_vector\_free(breaks);}
\NormalTok{\}}
\end{Highlighting}
\end{Shaded}

\hypertarget{sinhaspline.c}{%
\subsubsection{sinhaSpline.c}\label{sinhaspline.c}}

\begin{Shaded}
\begin{Highlighting}[]
\PreprocessorTok{\#include }\ImportTok{\textless{}stddef.h\textgreater{}}
\PreprocessorTok{\#include }\ImportTok{\textless{}stdio.h\textgreater{}}
\PreprocessorTok{\#include }\ImportTok{\textless{}stdlib.h\textgreater{}}

\DataTypeTok{double}\NormalTok{ bs(}\DataTypeTok{int}\NormalTok{ nknots, }\DataTypeTok{int}\NormalTok{ nspline, }\DataTypeTok{int}\NormalTok{ degree, }\DataTypeTok{double}\NormalTok{ x, }\DataTypeTok{double}\NormalTok{ *knots);}
\DataTypeTok{int}\NormalTok{ mindex(}\DataTypeTok{int}\NormalTok{ i, }\DataTypeTok{int}\NormalTok{ j, }\DataTypeTok{int}\NormalTok{ nrow);}

\DataTypeTok{void}\NormalTok{ sinhaBasis(}\DataTypeTok{int}\NormalTok{ *d, }\DataTypeTok{int}\NormalTok{ *n, }\DataTypeTok{int}\NormalTok{ *m, }\DataTypeTok{double}\NormalTok{ *x, }\DataTypeTok{double}\NormalTok{ *knots,}
                \DataTypeTok{double}\NormalTok{ *basis) \{}
    \DataTypeTok{int}\NormalTok{ mm = *m, dd = *d, nn = *n;}
    \DataTypeTok{int}\NormalTok{ k = mm {-} dd {-} }\DecValTok{1}\NormalTok{, i, j, ir, jr;}
    \ControlFlowTok{for}\NormalTok{ (i = }\DecValTok{0}\NormalTok{; i \textless{} nn; i++) \{}
\NormalTok{        ir = i + }\DecValTok{1}\NormalTok{;}
        \ControlFlowTok{if}\NormalTok{ (x[i] == knots[mm {-} }\DecValTok{1}\NormalTok{]) \{}
\NormalTok{            basis[mindex(ir, k, nn) {-} }\DecValTok{1}\NormalTok{] = }\FloatTok{1.0}\NormalTok{;}
            \ControlFlowTok{for}\NormalTok{ (j = }\DecValTok{0}\NormalTok{; j \textless{} (k {-} }\DecValTok{1}\NormalTok{); j++) \{}
\NormalTok{                jr = j + }\DecValTok{1}\NormalTok{;}
\NormalTok{                basis[mindex(ir, jr, nn) {-} }\DecValTok{1}\NormalTok{] = }\FloatTok{0.0}\NormalTok{;}
\NormalTok{            \}}
\NormalTok{        \} }\ControlFlowTok{else}\NormalTok{ \{}
            \ControlFlowTok{for}\NormalTok{ (j = }\DecValTok{0}\NormalTok{; j \textless{} k; j++) \{}
\NormalTok{                jr = j + }\DecValTok{1}\NormalTok{;}
\NormalTok{                basis[mindex(ir, jr, nn) {-} }\DecValTok{1}\NormalTok{] = bs(mm, jr, dd + }\DecValTok{1}\NormalTok{, x[i], knots);}
\NormalTok{            \}}
\NormalTok{        \}}
\NormalTok{    \}}
\NormalTok{\}}

\DataTypeTok{int}\NormalTok{ mindex(}\DataTypeTok{int}\NormalTok{ i, }\DataTypeTok{int}\NormalTok{ j, }\DataTypeTok{int}\NormalTok{ nrow) \{ }\ControlFlowTok{return}\NormalTok{ (j {-} }\DecValTok{1}\NormalTok{) * nrow + i; \}}

\DataTypeTok{double}\NormalTok{ bs(}\DataTypeTok{int}\NormalTok{ nknots, }\DataTypeTok{int}\NormalTok{ nspline, }\DataTypeTok{int}\NormalTok{ updegree, }\DataTypeTok{double}\NormalTok{ x, }\DataTypeTok{double}\NormalTok{ *knots) \{}
    \DataTypeTok{double}\NormalTok{ y, y1, y2, temp1, temp2;}
    \ControlFlowTok{if}\NormalTok{ (updegree == }\DecValTok{1}\NormalTok{) \{}
        \ControlFlowTok{if}\NormalTok{ ((x \textgreater{}= knots[nspline {-} }\DecValTok{1}\NormalTok{]) \&\& (x \textless{} knots[nspline]))}
\NormalTok{            y = }\FloatTok{1.0}\NormalTok{;}
        \ControlFlowTok{else}
\NormalTok{            y = }\FloatTok{0.0}\NormalTok{;}
\NormalTok{    \} }\ControlFlowTok{else}\NormalTok{ \{}
\NormalTok{        temp1 = }\FloatTok{0.0}\NormalTok{;}
        \ControlFlowTok{if}\NormalTok{ ((knots[nspline + updegree {-} }\DecValTok{2}\NormalTok{] {-} knots[nspline {-} }\DecValTok{1}\NormalTok{]) \textgreater{} }\DecValTok{0}\NormalTok{)}
\NormalTok{            temp1 = (x {-} knots[nspline {-} }\DecValTok{1}\NormalTok{]) /}
\NormalTok{                    (knots[nspline + updegree {-} }\DecValTok{2}\NormalTok{] {-} knots[nspline {-} }\DecValTok{1}\NormalTok{]);}
\NormalTok{        temp2 = }\FloatTok{0.0}\NormalTok{;}
        \ControlFlowTok{if}\NormalTok{ ((knots[nspline + updegree {-} }\DecValTok{1}\NormalTok{] {-} knots[nspline]) \textgreater{} }\DecValTok{0}\NormalTok{)}
\NormalTok{            temp2 = (knots[nspline + updegree {-} }\DecValTok{1}\NormalTok{] {-} x) /}
\NormalTok{                    (knots[nspline + updegree {-} }\DecValTok{1}\NormalTok{] {-} knots[nspline]);}
\NormalTok{        y1 = bs(nknots, nspline, updegree {-} }\DecValTok{1}\NormalTok{, x, knots);}
\NormalTok{        y2 = bs(nknots, nspline + }\DecValTok{1}\NormalTok{, updegree {-} }\DecValTok{1}\NormalTok{, x, knots);}
\NormalTok{        y = temp1 * y1 + temp2 * y2;}
\NormalTok{    \}}
    \ControlFlowTok{return}\NormalTok{ y;}
\NormalTok{\}}
\end{Highlighting}
\end{Shaded}

\hypertarget{deboor.c}{%
\subsubsection{deboor.c}\label{deboor.c}}

\begin{Shaded}
\begin{Highlighting}[]

\PreprocessorTok{\#include }\ImportTok{\textless{}math.h\textgreater{}}
\PreprocessorTok{\#include }\ImportTok{\textless{}stdbool.h\textgreater{}}
\PreprocessorTok{\#include }\ImportTok{\textless{}stdlib.h\textgreater{}}

\KeywordTok{inline} \DataTypeTok{int}\NormalTok{ VINDEX(}\DataTypeTok{const} \DataTypeTok{int}\NormalTok{);}
\KeywordTok{inline} \DataTypeTok{int}\NormalTok{ MINDEX(}\DataTypeTok{const} \DataTypeTok{int}\NormalTok{, }\DataTypeTok{const} \DataTypeTok{int}\NormalTok{, }\DataTypeTok{const} \DataTypeTok{int}\NormalTok{);}

\DataTypeTok{void}\NormalTok{ checkIncreasing(}\DataTypeTok{const} \DataTypeTok{double}\NormalTok{ *, }\DataTypeTok{const} \DataTypeTok{double}\NormalTok{ *, }\DataTypeTok{const} \DataTypeTok{double}\NormalTok{ *,}
                     \DataTypeTok{const} \DataTypeTok{int}\NormalTok{ *, }\DataTypeTok{bool}\NormalTok{ *);}
\DataTypeTok{void}\NormalTok{ extendPartition(}\DataTypeTok{const} \DataTypeTok{double}\NormalTok{ *, }\DataTypeTok{const} \DataTypeTok{int}\NormalTok{ *, }\DataTypeTok{const} \DataTypeTok{int}\NormalTok{ *, }\DataTypeTok{const} \DataTypeTok{int}\NormalTok{ *,}
                     \DataTypeTok{const} \DataTypeTok{double}\NormalTok{ *, }\DataTypeTok{const} \DataTypeTok{double}\NormalTok{ *, }\DataTypeTok{double}\NormalTok{ *);}
\DataTypeTok{void}\NormalTok{ bisect(}\DataTypeTok{const} \DataTypeTok{double}\NormalTok{ *, }\DataTypeTok{const} \DataTypeTok{double}\NormalTok{ *, }\DataTypeTok{const} \DataTypeTok{int}\NormalTok{ *, }\DataTypeTok{const} \DataTypeTok{int}\NormalTok{ *, }\DataTypeTok{int}\NormalTok{ *);}
\DataTypeTok{void}\NormalTok{ bsplines(}\DataTypeTok{const} \DataTypeTok{double}\NormalTok{ *, }\DataTypeTok{const} \DataTypeTok{double}\NormalTok{ *, }\DataTypeTok{const} \DataTypeTok{int}\NormalTok{ *, }\DataTypeTok{const} \DataTypeTok{int}\NormalTok{ *, }\DataTypeTok{int}\NormalTok{ *,}
              \DataTypeTok{double}\NormalTok{ *);}
\DataTypeTok{void}\NormalTok{ bsplineBasis(}\DataTypeTok{const} \DataTypeTok{double}\NormalTok{ *, }\DataTypeTok{const} \DataTypeTok{double}\NormalTok{ *, }\DataTypeTok{const} \DataTypeTok{int}\NormalTok{ *, }\DataTypeTok{const} \DataTypeTok{int}\NormalTok{ *,}
                  \DataTypeTok{const} \DataTypeTok{int}\NormalTok{ *, }\DataTypeTok{double}\NormalTok{ *);}
\DataTypeTok{void}\NormalTok{ isplineBasis(}\DataTypeTok{const} \DataTypeTok{double}\NormalTok{ *, }\DataTypeTok{const} \DataTypeTok{double}\NormalTok{ *, }\DataTypeTok{const} \DataTypeTok{int}\NormalTok{ *, }\DataTypeTok{const} \DataTypeTok{int}\NormalTok{ *,}
                  \DataTypeTok{const} \DataTypeTok{int}\NormalTok{ *, }\DataTypeTok{double}\NormalTok{ *);}
\DataTypeTok{void}\NormalTok{ bsplineSparse(}\DataTypeTok{const} \DataTypeTok{double}\NormalTok{ *, }\DataTypeTok{const} \DataTypeTok{double}\NormalTok{ *, }\DataTypeTok{const} \DataTypeTok{int}\NormalTok{ *, }\DataTypeTok{const} \DataTypeTok{int}\NormalTok{ *,}
                   \DataTypeTok{const} \DataTypeTok{int}\NormalTok{ *, }\DataTypeTok{int}\NormalTok{ *, }\DataTypeTok{double}\NormalTok{ *);}
\DataTypeTok{void}\NormalTok{ isplineSparse(}\DataTypeTok{const} \DataTypeTok{double}\NormalTok{ *, }\DataTypeTok{const} \DataTypeTok{double}\NormalTok{ *, }\DataTypeTok{const} \DataTypeTok{int}\NormalTok{ *, }\DataTypeTok{const} \DataTypeTok{int}\NormalTok{ *,}
                   \DataTypeTok{const} \DataTypeTok{int}\NormalTok{ *, }\DataTypeTok{int}\NormalTok{ *, }\DataTypeTok{double}\NormalTok{ *);}

\KeywordTok{inline} \DataTypeTok{int}\NormalTok{ VINDEX(}\DataTypeTok{const} \DataTypeTok{int}\NormalTok{ i) \{ }\ControlFlowTok{return}\NormalTok{ i {-} }\DecValTok{1}\NormalTok{; \}}

\KeywordTok{inline} \DataTypeTok{int}\NormalTok{ MINDEX(}\DataTypeTok{const} \DataTypeTok{int}\NormalTok{ i, }\DataTypeTok{const} \DataTypeTok{int}\NormalTok{ j, }\DataTypeTok{const} \DataTypeTok{int}\NormalTok{ n) \{}
    \ControlFlowTok{return}\NormalTok{ (i {-} }\DecValTok{1}\NormalTok{) + (j {-} }\DecValTok{1}\NormalTok{) * n;}
\NormalTok{\}}

\KeywordTok{inline} \DataTypeTok{int}\NormalTok{ IMIN(}\DataTypeTok{const} \DataTypeTok{int}\NormalTok{ a, }\DataTypeTok{const} \DataTypeTok{int}\NormalTok{ b) \{}
    \ControlFlowTok{if}\NormalTok{ (a \textgreater{} b) }\ControlFlowTok{return}\NormalTok{ b;}
    \ControlFlowTok{return}\NormalTok{ a;}
\NormalTok{\}}

\KeywordTok{inline} \DataTypeTok{int}\NormalTok{ IMAX(}\DataTypeTok{const} \DataTypeTok{int}\NormalTok{ a, }\DataTypeTok{const} \DataTypeTok{int}\NormalTok{ b) \{}
    \ControlFlowTok{if}\NormalTok{ (a \textless{} b) }\ControlFlowTok{return}\NormalTok{ b;}
    \ControlFlowTok{return}\NormalTok{ a;}
\NormalTok{\}}

\DataTypeTok{void}\NormalTok{ checkIncreasing(}\DataTypeTok{const} \DataTypeTok{double}\NormalTok{ *innerknots, }\DataTypeTok{const} \DataTypeTok{double}\NormalTok{ *lowend,}
                     \DataTypeTok{const} \DataTypeTok{double}\NormalTok{ *highend, }\DataTypeTok{const} \DataTypeTok{int}\NormalTok{ *ninner, }\DataTypeTok{bool}\NormalTok{ *fail) \{}
\NormalTok{    *fail = false;}
    \ControlFlowTok{if}\NormalTok{ (*lowend \textgreater{}= innerknots[VINDEX(}\DecValTok{1}\NormalTok{)]) \{}
\NormalTok{        *fail = true;}
        \ControlFlowTok{return}\NormalTok{;}
\NormalTok{    \}}
    \ControlFlowTok{if}\NormalTok{ (*highend \textless{}= innerknots[VINDEX(*ninner)]) \{}
\NormalTok{        *fail = true;}
        \ControlFlowTok{return}\NormalTok{;}
\NormalTok{    \}}
    \ControlFlowTok{for}\NormalTok{ (}\DataTypeTok{int}\NormalTok{ i = }\DecValTok{1}\NormalTok{; i \textless{} *ninner; i++) \{}
        \ControlFlowTok{if}\NormalTok{ (innerknots[i] \textless{}= innerknots[i {-} }\DecValTok{1}\NormalTok{]) \{}
\NormalTok{            *fail = true;}
            \ControlFlowTok{return}\NormalTok{;}
\NormalTok{        \}}
\NormalTok{    \}}
\NormalTok{\}}

\DataTypeTok{void}\NormalTok{ extendPartition(}\DataTypeTok{const} \DataTypeTok{double}\NormalTok{ *innerknots, }\DataTypeTok{const} \DataTypeTok{int}\NormalTok{ *multiplicities,}
                     \DataTypeTok{const} \DataTypeTok{int}\NormalTok{ *order, }\DataTypeTok{const} \DataTypeTok{int}\NormalTok{ *ninner, }\DataTypeTok{const} \DataTypeTok{double}\NormalTok{ *lowend,}
                     \DataTypeTok{const} \DataTypeTok{double}\NormalTok{ *highend, }\DataTypeTok{double}\NormalTok{ *extended) \{}
    \DataTypeTok{int}\NormalTok{ k = }\DecValTok{1}\NormalTok{;}
    \ControlFlowTok{for}\NormalTok{ (}\DataTypeTok{int}\NormalTok{ i = }\DecValTok{1}\NormalTok{; i \textless{}= *order; i++) \{}
\NormalTok{        extended[VINDEX(k)] = *lowend;}
\NormalTok{        k++;}
\NormalTok{    \}}
    \ControlFlowTok{for}\NormalTok{ (}\DataTypeTok{int}\NormalTok{ j = }\DecValTok{1}\NormalTok{; j \textless{}= *ninner; j++)}
        \ControlFlowTok{for}\NormalTok{ (}\DataTypeTok{int}\NormalTok{ i = }\DecValTok{1}\NormalTok{; i \textless{}= multiplicities[VINDEX(j)]; i++) \{}
\NormalTok{            extended[VINDEX(k)] = innerknots[VINDEX(j)];}
\NormalTok{            k++;}
\NormalTok{        \}}
    \ControlFlowTok{for}\NormalTok{ (}\DataTypeTok{int}\NormalTok{ i = }\DecValTok{1}\NormalTok{; i \textless{}= *order; i++) \{}
\NormalTok{        extended[VINDEX(k)] = *highend;}
\NormalTok{        k++;}
\NormalTok{    \}}
\NormalTok{\}}

\DataTypeTok{void}\NormalTok{ bisect(}\DataTypeTok{const} \DataTypeTok{double}\NormalTok{ *x, }\DataTypeTok{const} \DataTypeTok{double}\NormalTok{ *knots, }\DataTypeTok{const} \DataTypeTok{int}\NormalTok{ *lowindex,}
            \DataTypeTok{const} \DataTypeTok{int}\NormalTok{ *highindex, }\DataTypeTok{int}\NormalTok{ *index) \{}
    \DataTypeTok{int}\NormalTok{ l = *lowindex, u = *highindex, mid = }\DecValTok{0}\NormalTok{;}
    \ControlFlowTok{while}\NormalTok{ ((u {-} l) \textgreater{} }\DecValTok{1}\NormalTok{) \{}
\NormalTok{        mid = (}\DataTypeTok{int}\NormalTok{)floor((u + l) / }\DecValTok{2}\NormalTok{);}
        \ControlFlowTok{if}\NormalTok{ (*x \textless{} knots[VINDEX(mid)])}
\NormalTok{            u = mid;}
        \ControlFlowTok{else}
\NormalTok{            l = mid;}
\NormalTok{    \}}
\NormalTok{    *index = l;}
    \ControlFlowTok{return}\NormalTok{;}
\NormalTok{\}}

\DataTypeTok{void}\NormalTok{ bsplines(}\DataTypeTok{const} \DataTypeTok{double}\NormalTok{ *x, }\DataTypeTok{const} \DataTypeTok{double}\NormalTok{ *knots, }\DataTypeTok{const} \DataTypeTok{int}\NormalTok{ *order,}
              \DataTypeTok{const} \DataTypeTok{int}\NormalTok{ *nknots, }\DataTypeTok{int}\NormalTok{ *index, }\DataTypeTok{double}\NormalTok{ *q) \{}
    \DataTypeTok{int}\NormalTok{ lowindex = }\DecValTok{1}\NormalTok{, highindex = *nknots, m = *order, j, jp1;}
    \DataTypeTok{double}\NormalTok{ drr, dll, saved, term;}
    \DataTypeTok{double}\NormalTok{ *dr = (}\DataTypeTok{double}\NormalTok{ *)calloc((}\DataTypeTok{size\_t}\NormalTok{)m, }\KeywordTok{sizeof}\NormalTok{(}\DataTypeTok{double}\NormalTok{));}
    \DataTypeTok{double}\NormalTok{ *dl = (}\DataTypeTok{double}\NormalTok{ *)calloc((}\DataTypeTok{size\_t}\NormalTok{)m, }\KeywordTok{sizeof}\NormalTok{(}\DataTypeTok{double}\NormalTok{));}
\NormalTok{    (}\DataTypeTok{void}\NormalTok{)bisect(x, knots, \&lowindex, \&highindex, index);}
    \DataTypeTok{int}\NormalTok{ l = *index;}
    \ControlFlowTok{for}\NormalTok{ (j = }\DecValTok{1}\NormalTok{; j \textless{}= m; j++) \{}
\NormalTok{        q[VINDEX(j)] = }\FloatTok{0.0}\NormalTok{;}
\NormalTok{    \}}
    \ControlFlowTok{if}\NormalTok{ (*x == knots[VINDEX(*nknots)]) \{}
\NormalTok{        q[VINDEX(m)] = }\FloatTok{1.0}\NormalTok{;}
        \ControlFlowTok{return}\NormalTok{;}
\NormalTok{    \}}
\NormalTok{    q[VINDEX(}\DecValTok{1}\NormalTok{)] = }\FloatTok{1.0}\NormalTok{;}
\NormalTok{    j = }\DecValTok{1}\NormalTok{;}
    \ControlFlowTok{if}\NormalTok{ (j \textgreater{}= m) }\ControlFlowTok{return}\NormalTok{;}
    \ControlFlowTok{while}\NormalTok{ (j \textless{} m) \{}
\NormalTok{        dr[VINDEX(j)] = knots[VINDEX(l + j)] {-} *x;}
\NormalTok{        dl[VINDEX(j)] = *x {-} knots[VINDEX(l + }\DecValTok{1}\NormalTok{ {-} j)];}
\NormalTok{        jp1 = j + }\DecValTok{1}\NormalTok{;}
\NormalTok{        saved = }\FloatTok{0.0}\NormalTok{;}
        \ControlFlowTok{for}\NormalTok{ (}\DataTypeTok{int}\NormalTok{ r = }\DecValTok{1}\NormalTok{; r \textless{}= j; r++) \{}
\NormalTok{            drr = dr[VINDEX(r)];}
\NormalTok{            dll = dl[VINDEX(jp1 {-} r)];}
\NormalTok{            term = q[VINDEX(r)] / (drr + dll);}
\NormalTok{            q[VINDEX(r)] = saved + drr * term;}
\NormalTok{            saved = dll * term;}
\NormalTok{        \}}
\NormalTok{        q[VINDEX(jp1)] = saved;}
\NormalTok{        j = jp1;}
\NormalTok{    \}}
\NormalTok{    free(dr);}
\NormalTok{    free(dl);}
    \ControlFlowTok{return}\NormalTok{;}
\NormalTok{\}}

\DataTypeTok{void}\NormalTok{ bsplineBasis(}\DataTypeTok{const} \DataTypeTok{double}\NormalTok{ *x, }\DataTypeTok{const} \DataTypeTok{double}\NormalTok{ *knots, }\DataTypeTok{const} \DataTypeTok{int}\NormalTok{ *order,}
                  \DataTypeTok{const} \DataTypeTok{int}\NormalTok{ *nknots, }\DataTypeTok{const} \DataTypeTok{int}\NormalTok{ *nvalues, }\DataTypeTok{double}\NormalTok{ *result) \{}
    \DataTypeTok{int}\NormalTok{ m = *order, l = }\DecValTok{0}\NormalTok{;}
    \DataTypeTok{double}\NormalTok{ *q = (}\DataTypeTok{double}\NormalTok{ *)calloc((}\DataTypeTok{size\_t}\NormalTok{)m + }\DecValTok{1}\NormalTok{, }\KeywordTok{sizeof}\NormalTok{(}\DataTypeTok{double}\NormalTok{));}
    \ControlFlowTok{for}\NormalTok{ (}\DataTypeTok{int}\NormalTok{ i = }\DecValTok{1}\NormalTok{; i \textless{}= *nvalues; i++) \{}
\NormalTok{        (}\DataTypeTok{void}\NormalTok{)bsplines(x + VINDEX(i), knots, order, nknots, \&l, q);}
        \ControlFlowTok{for}\NormalTok{ (}\DataTypeTok{int}\NormalTok{ j = }\DecValTok{1}\NormalTok{; j \textless{}= m; j++) \{}
            \DataTypeTok{int}\NormalTok{ r = IMIN(l {-} m + j, *nknots {-} m);}
\NormalTok{            result[MINDEX(i, r, *nvalues)] = q[VINDEX(j)];}
\NormalTok{        \}}
\NormalTok{    \}}
\NormalTok{    free(q);}
    \ControlFlowTok{return}\NormalTok{;}
\NormalTok{\}}

\DataTypeTok{void}\NormalTok{ isplineBasis(}\DataTypeTok{const} \DataTypeTok{double}\NormalTok{ *x, }\DataTypeTok{const} \DataTypeTok{double}\NormalTok{ *knots, }\DataTypeTok{const} \DataTypeTok{int}\NormalTok{ *order,}
                  \DataTypeTok{const} \DataTypeTok{int}\NormalTok{ *nknots, }\DataTypeTok{const} \DataTypeTok{int}\NormalTok{ *nvalues, }\DataTypeTok{double}\NormalTok{ *result) \{}
    \DataTypeTok{int}\NormalTok{ m = *nknots {-} *order, n = *nvalues;}
\NormalTok{    (}\DataTypeTok{void}\NormalTok{)bsplineBasis(x, knots, order, nknots, nvalues, result);}
    \ControlFlowTok{for}\NormalTok{ (}\DataTypeTok{int}\NormalTok{ i = }\DecValTok{1}\NormalTok{; i \textless{}= n; i++) \{}
        \ControlFlowTok{for}\NormalTok{ (}\DataTypeTok{int}\NormalTok{ j = m {-} }\DecValTok{1}\NormalTok{; j \textgreater{}= }\DecValTok{1}\NormalTok{; j{-}{-}) \{}
\NormalTok{            result[MINDEX(i, j, n)] += result[MINDEX(i, j + }\DecValTok{1}\NormalTok{, n)];}
\NormalTok{        \}}
\NormalTok{    \}}
    \ControlFlowTok{return}\NormalTok{;}
\NormalTok{\}}

\DataTypeTok{void}\NormalTok{ bsplineSparse(}\DataTypeTok{const} \DataTypeTok{double}\NormalTok{ *x, }\DataTypeTok{const} \DataTypeTok{double}\NormalTok{ *knots, }\DataTypeTok{const} \DataTypeTok{int}\NormalTok{ *order,}
                   \DataTypeTok{const} \DataTypeTok{int}\NormalTok{ *nknots, }\DataTypeTok{const} \DataTypeTok{int}\NormalTok{ *nvalues, }\DataTypeTok{int}\NormalTok{ *columns,}
                   \DataTypeTok{double}\NormalTok{ *result) \{}
    \DataTypeTok{int}\NormalTok{ m = *order, l = }\DecValTok{0}\NormalTok{;}
    \DataTypeTok{double}\NormalTok{ *q = (}\DataTypeTok{double}\NormalTok{ *)calloc((}\DataTypeTok{size\_t}\NormalTok{)m + }\DecValTok{1}\NormalTok{, }\KeywordTok{sizeof}\NormalTok{(}\DataTypeTok{double}\NormalTok{));}
    \ControlFlowTok{for}\NormalTok{ (}\DataTypeTok{int}\NormalTok{ i = }\DecValTok{1}\NormalTok{; i \textless{}= *nvalues; i++) \{}
\NormalTok{        (}\DataTypeTok{void}\NormalTok{)bsplines(x + VINDEX(i), knots, order, nknots, \&l, q);}
\NormalTok{        columns[VINDEX(i)] = l;}
        \ControlFlowTok{for}\NormalTok{ (}\DataTypeTok{int}\NormalTok{ j = }\DecValTok{1}\NormalTok{; j \textless{}= m; j++) \{}
\NormalTok{            result[MINDEX(i, j, *nvalues)] = q[VINDEX(j)];}
\NormalTok{        \}}
\NormalTok{    \}}
\NormalTok{    free(q);}
    \ControlFlowTok{return}\NormalTok{;}
\NormalTok{\}}

\DataTypeTok{void}\NormalTok{ isplineSparse(}\DataTypeTok{const} \DataTypeTok{double}\NormalTok{ *x, }\DataTypeTok{const} \DataTypeTok{double}\NormalTok{ *knots, }\DataTypeTok{const} \DataTypeTok{int}\NormalTok{ *order,}
                   \DataTypeTok{const} \DataTypeTok{int}\NormalTok{ *nknots, }\DataTypeTok{const} \DataTypeTok{int}\NormalTok{ *nvalues, }\DataTypeTok{int}\NormalTok{ *columns,}
                   \DataTypeTok{double}\NormalTok{ *result) \{}
    \DataTypeTok{int}\NormalTok{ m = *nknots {-} *order, n = *nvalues;}
    \ControlFlowTok{for}\NormalTok{ (}\DataTypeTok{int}\NormalTok{ i = }\DecValTok{1}\NormalTok{; i \textless{}= *nvalues; i++) \{}
\NormalTok{        (}\DataTypeTok{void}\NormalTok{)bsplineSparse(x + VINDEX(i), knots, order, nknots, nvalues,}
\NormalTok{                            columns, result);}
        \ControlFlowTok{for}\NormalTok{ (}\DataTypeTok{int}\NormalTok{ j = m {-} }\DecValTok{1}\NormalTok{; j \textgreater{}= }\DecValTok{1}\NormalTok{; j{-}{-}) \{}
\NormalTok{            result[MINDEX(i, j, n)] += result[MINDEX(i, j + }\DecValTok{1}\NormalTok{, n)];}
\NormalTok{        \}}
\NormalTok{    \}}
    \ControlFlowTok{return}\NormalTok{;}
\NormalTok{\}}
\end{Highlighting}
\end{Shaded}

\hypertarget{references}{%
\section*{References}\label{references}}
\addcontentsline{toc}{section}{References}

\hypertarget{refs}{}
\begin{cslreferences}
\leavevmode\hypertarget{ref-cox_72}{}%
Cox, M. G. 1972. ``The Numerical Evaluation of B-splines.''
\emph{Journal of the Institute of Mathematics and Its Applications} 10:
134--49.

\leavevmode\hypertarget{ref-deboor_72}{}%
De Boor, C. 1972. ``On Calculating with B-splines. II. Integration.''
\emph{Journal of Approximation Theory} 6: 50--62.

\leavevmode\hypertarget{ref-deboor_76}{}%
---------. 1976. ``Splines as Linear Combination of B-splines. A
Survey.'' In \emph{Approximation Theory II}, edited by G. G. Lorentz, C.
K. Chui, and L. L. Schumaker, 1--47. Academic Press.

\leavevmode\hypertarget{ref-deboor_01}{}%
---------. 2001. \emph{A Practical Guide to Splines}. Revised Edition.
New York: Springer-Verlag.

\leavevmode\hypertarget{ref-deboor_hollig_85}{}%
De Boor, C., and K. Höllig. 1985. ``B-splines without Divided
Differences.'' Technical Report 622. Department of Computer Science,
University of Wisconsin-Madison.

\leavevmode\hypertarget{ref-deboor_lyche_schumaker_76}{}%
De Boor, C., T. Lyche, and L. L. Schumaker. 1976. ``On Calculating with
B-splines. II. Integration.'' In \emph{Numerische Methoden der
Approximationstheorie}, edited by L. Collatz, G. Meinardus, and H.
Werner, 123--46. Basel: Birkhauser.

\leavevmode\hypertarget{ref-deleeuw_E_15d}{}%
De Leeuw, J. 2015. ``Exceedingly Simple B-Spline Code.''
\url{http://rpubs.com/deleeuw/79161}.

\leavevmode\hypertarget{ref-deleeuw_E_15e}{}%
De Leeuw, J. 2015. ``Regression with Linear Inequality Restrictions on
Predicted Values.'' \url{http://rpubs.com/deleeuw/78897}.

\leavevmode\hypertarget{ref-deleeuw_mair_A_09a}{}%
De Leeuw, J., and P. Mair. 2009. ``Homogeneity Analysis in R: The
Package Homals.'' \emph{Journal of Statistical Software} 31 (4): 1--21.
\url{http://deleeuwpdx.net/janspubs/2009/articles/deleeuw_mair_A_09a.pdf}.

\leavevmode\hypertarget{ref-gaffney_76}{}%
Gaffney, P. W. 1976. ``The Calculation of Indefinite Integrals of
B-splines.'' \emph{Journal of the Institute of Mathematics and Its
Applications} 17: 37--41.

\leavevmode\hypertarget{ref-gsl_manual_16}{}%
Galassi, M., J. Davies, J. Theiler, B. Gough, G. Jungman, P. Alken, M.
Booth, F. Rossi, and R. Ulerich. 2016. ``GNU Scienfific Library,
Reference Manual.'' In, Edition 2.3.
\url{https://www.gnu.org/software/gsl/manual/html_node/Basis-Splines.html}.

\leavevmode\hypertarget{ref-gifi_B_90}{}%
Gifi, A. 1990. \emph{Nonlinear Multivariate Analysis}. New York, N.Y.:
Wiley.

\leavevmode\hypertarget{ref-ramsay_88}{}%
Ramsay, J. O. 1988. ``Monotone Regression Splines in Action.''
\emph{Statistical Science} 3: 425--61.

\leavevmode\hypertarget{ref-schumaker_07}{}%
Schumaker, L. 2007. \emph{Spline Functions: Basic Theory}. Third
Edition. Cambridge University Press.

\leavevmode\hypertarget{ref-sinha}{}%
Sinha, S. n.d. ``A Very Short Note on B-splines.''
\url{https://www.stat.tamu.edu/~sinha/research/note1.pdf}.

\leavevmode\hypertarget{ref-wang_lawson_hanson_15}{}%
Wang, Y., C. L. Lawson, and R. J. Hanson. 2015. \emph{lsei: Solving
Least Squares Problems under Equality/Inequality Constraints}.
\url{http://CRAN.R-project.org/package=lsei}.
\end{cslreferences}

\end{document}
